\documentclass[11pt]{article}
\usepackage[localtheorems]{raeez-math-template}

%%% Packages

%%% Page geometry

%%% Theorem environments
\theoremstyle{plain}
\newtheorem{theorem}{Theorem}[section]
\newtheorem{proposition}[theorem]{Proposition}
\newtheorem{corollary}[theorem]{Corollary}
\newtheorem{lemma}[theorem]{Lemma}
\newtheorem{conjecture}[theorem]{Conjecture}

\theoremstyle{definition}
\newtheorem{definition}[theorem]{Definition}
\newtheorem{construction}[theorem]{Construction}
\newtheorem{example}[theorem]{Example}

\theoremstyle{remark}
\newtheorem{remark}[theorem]{Remark}

%%% Macros
\newcommand{\fg}{\mathfrak{g}}
\newcommand{\fh}{\mathfrak{h}}
\newcommand{\fsl}{\mathfrak{sl}}
\newcommand{\Ran}{\mathrm{Ran}}
\newcommand{\Mbar}{\overline{\mathcal{M}}}
\newcommand{\FM}{\mathrm{FM}}
\newcommand{\MC}{\mathrm{MC}}
\newcommand{\cA}{\mathcal{A}}
\newcommand{\cC}{\mathcal{C}}
\newcommand{\cH}{\mathcal{H}}
\newcommand{\cW}{\mathcal{W}}
\newcommand{\cZ}{\mathcal{Z}}
\newcommand{\cD}{\mathcal{D}}
\newcommand{\cM}{\mathcal{M}}
\newcommand{\barB}{\overline{B}}
\newcommand{\Abar}{\bar{\cA}}
\newcommand{\Res}{\operatorname{Res}}
\newcommand{\disc}{\operatorname{disc}}
\newcommand{\ChirHoch}{\mathrm{ChirHoch}}
\newcommand{\Eone}{E_1}
\newcommand{\Einf}{E_\infty}
\newcommand{\gAmod}{\mathfrak{g}_{\cA}^{\mathrm{mod}}}
\newcommand{\Ahat}{\hat{A}}
\newcommand{\Confinf}{\mathrm{Conv}_\infty}
\DeclareMathOperator{\Aut}{Aut}
\DeclareMathOperator{\Hom}{Hom}
\DeclareMathOperator{\Conf}{Conf}
\DeclareMathOperator{\obs}{obs}

%%% Title
\title{The five theorems of modular Koszul duality}
\author{Raeez Lorgat}
\date{}

\begin{document}
\maketitle

%%% ================================================================
%%% ABSTRACT
%%% ================================================================

\begin{abstract}
Koszul duality over a point closes through a quadratic coalgebra, the
full bar construction, and cobar reconstruction.  Over an algebraic
curve, Ran collisions and the geometry of moduli add further layers.
The enhanced bar differential in $D(\Ran X)$ is square-zero.  After
second-kind cobar realization, its fiberwise modular algebra is curved,
$m_1^2=[m_0,-]$.  A normalized collision trace extracts the scalar
$\kappa(\cA)$, while the full curvature retains operator-valued data.

The chiral algebra $\cA$ is the level-$1$ chart
$A_b = \mathrm{End}_{\mathcal{C}}(b)$ of a factorization
dg-category $\mathcal{C}^{\mathrm{op}}$ on a tangential log
curve $(X, D, \tau)$, evaluated at a vacuum object $b$. Five
theorems characterize the levels of the resulting Beilinson
diagram
\[
\begin{tikzcd}[column sep=large,row sep=small]
 \mathcal C^{\mathrm{op}}
 \arrow[r,"b"]
 & A_b
 \arrow[r,"\bar B"]
 \arrow[d,"Z^{\mathrm{der}}_{\mathrm{ch}}"']
 & B_X(A_b)
 \arrow[d,"\mathbb D_{\Ran}"] \\
 & C^\bullet_{\mathrm{ch}}(A_b,A_b)
 & K_X(A_b).
\end{tikzcd}
\]
Here
\(C^\bullet_{\mathrm{ch}}(A_b,A_b)
:=R\!\operatorname{Hom}_{A_b\otimes A_b^{\mathrm{op}}}(A_b,A_b)\)
is the algebraic derived centre in the completed chiral bimodule
category.  A physical bulk realization consists of an open--closed
equivalence
\[
 \Psi_{\mathrm{OCA}}\colon
 \mathcal O^{\mathrm{phys}}_{\mathrm{bulk}}(A_b)
 \xrightarrow{\ \sim\ }C^\bullet_{\mathrm{ch}}(A_b,A_b).
\]
The line-operator and modular-trace stages carry their respective
comparison data.
Each theorem carries a type signature
(\emph{quadrant, presentation, level, hypothesis package}).

Theorem~A separates Ran reconstruction
$R_X(\cA)=\Omega_XB_X(\cA)\xrightarrow{\sim}\cA$ from the
Verdier--Koszul algebra
$K_X(\cA)=\mathbb D_{\Ran}B_X(\cA)$.  The ambient Ran equivalence is
the theorem of Francis--Gaitsgory; its restriction to the chosen
factorization subcategories carries the package $H_{\mathrm{fact}}$.
Theorem~B is the quadratic recognition theorem: under
$H_{\mathrm{CL}}$, the map
$q_\cA\colon\cA^{\mathrm i}\to B_X(\cA)$ is a quasi-isomorphism
precisely when the associated twisted tensor products are acyclic
and the quadratic bar spectral sequence converges with diagonal
support.
Theorem~C is the $5 \times 5$ $\kappa$-stratification matrix
on rows $\mathsf{G/L/C/M/B}$ with columns
$\{\kappa_{\mathrm{cat}},
\kappa^{\mathrm{Hodge}}_{\mathrm{ch}},
\kappa^{\mathrm{Heis}}_{\mathrm{ch}},
\kappa_{\mathrm{BKM}},
\kappa_{\mathrm{fiber}}\}$ at level $5$ per stratum;
$K3 \times E$ has native entries $(0,\,0,\,\star,\,5,\,24)$, and its
Open Heisenberg entry~$3$ is conditional on the realization map
\(\eta_{K3,E}\).  On a family carrying the represented
anti-invariant pairing in $H_C$, the Lagrangian decomposition yields
the corresponding duality-orbit sum; the standard Kac--Moody,
free-field, and Virasoro traces give $0$, $0$, and $13$ on their
respective comparison surfaces.
Theorem~D separates four typed objects:
$\operatorname{Obs}^{\mathrm{def}}_g\in H^2(\mathrm{Def}_g)$ for
$g\ge2$,
$\mathfrak O_g^K=\kappa\lambda_{-1}(\mathbb E_g)$ under $H_D^K$ for
$g\ge2$,
$\operatorname{obs}^{\mathrm{Hdg}}_g=(-1)^g\kappa\lambda_g$ for
$g\ge2$, and
$F_g=\kappa\lambda_g^{\mathrm{FP}}+\delta F_g^{\mathrm{cross}}$
under the stable-graph package.  The genus-one trace has the
normalization
$\operatorname{tr}_1\operatorname{Obs}^{\mathrm{def}}_{1,1}
=\kappa\lambda_1$.
Theorem~H is family-indexed: a strong deformation retract
$H_H(A_b;S)$ identifies the chiral Hochschild complex with a model
supported in~$S$.  The bounded Virasoro and rank-one even superboson
calculations enter the chart complex through a specified comparison
quasi-isomorphism.

The ordered bar complex
$\barB^{\mathrm{ord}}_X(A_b)$ with its deconcatenation coproduct and
collision kernel $K_{A_b}^{\mathrm{coll}}(z)$ is a common source for
the shadow operators.  Their scalar coordinates use normalized
equivariant functionals.  In arity two,
\(q_2(K_\cA^{\mathrm{coll}})\in
\operatorname{End}_\cA(2)_{h\Sigma_2}\), and
\(\tau_{\cA,2}^{\mathrm{res}}\colon
\operatorname{End}_\cA(2)_{h\Sigma_2}\to\mathbf C\)
recovers $\kappa$ on the abelian and scalar lanes; the non-abelian
Kac--Moody value
\(\kappa=\tau_{\cA,2}^{\mathrm{res}}
q_2(K_\cA^{\mathrm{coll}})+\dim(\fg)/2\)
encodes the Sugawara shift. The locus
$\mathrm{KSDual} \subset \mathrm{Kosz}(X)$ on which the Verdier
involution $\cA \mapsto \cA^!$ acts by the identity is the
$\mathbb{Z}/2$-fixed sublocus in the finite-type curved
Verdier--Koszul ambient. On its admissible Koszul-complete part, a
chosen self-duality map
$\sigma_\cA\colon\cA\xrightarrow{\sim}K_X(\cA)$ represents the fixed
point, and Theorem~H retains its family support set~$S$.  The
Borcherds trace identity \cite{Borcherds1998,Gritsenko1999}
$\kappa_{\mathrm{BKM}}(\Phi_N) = c_N(0)/2$
is the equivariant signature.
The Heisenberg algebra provides the first exact local calculation and
displays the comparison maps required for its global realization.
\end{abstract}

\tableofcontents

%%% ================================================================
%%% SECTION 1: INTRODUCTION
%%% ================================================================

\section{Introduction: from a point to a curve}
\label{sec:introduction}

\subsection{The problem}

Koszul duality over a point begins with the full coalgebra
$B(A)=(T^c(s^{-1}\bar A),d_B)$ and the reconstruction counit
$\Omega B(A)\to A$ \cite{Priddy70,BGS96,LV12}.  A connected
quadratic presentation supplies $A^{\mathrm i}\to B(A)$; the Koszul
criterion identifies this map as a quasi-isomorphism and forms
$A^!=(A^{\mathrm i})^\vee$.

Over an algebraic curve, the generators are sections of a
$\cD_X$-module, and the relations are operator product expansions
with arbitrary pole order.  Higher-genus geometry places the
fiberwise cobar algebra in a curved $A_\infty$/CDG package:
the raw identity is $m_1^2=[m_0^{(g)},-]$, and its scalar diagonal
projection is $\kappa(\cA)\cdot\omega_g$.

The question is: what structure replaces the terminating classical
theory?

\subsection{The answer}

The chiral algebra $\cA = A_b$ is the chart at a vacuum
$b \in \mathcal{C}^{\mathrm{op}}$ of a factorization dg-category
on a tangential log curve $(X, D, \tau)$.  Its two algebraic branches
are the Ran bar--Verdier diagram
\(A_b\xrightarrow{\bar B}B_X(A_b)
\xrightarrow{\mathbb D_{\Ran}}K_X(A_b)\)
and the derived-centre construction
\(A_b\mapsto C^\bullet_{\mathrm{ch}}(A_b,A_b)\).
These branches organise levels~$1$--$3$ of the open quadrant.  The
open--closed equivalence \(\Psi_{\mathrm{OCA}}\) supplies a physical
bulk realization.
On the strictly augmented lane, the ordered bar complex
\begin{equation}\label{eq:bar-complex-def}
\barB^{\mathrm{ord}}_X(\cA)
\;=\;
T^c(s^{-1}\Abar),
\qquad
|s^{-1}v| = |v| - 1,
\end{equation}
with deconcatenation coproduct
\begin{equation}\label{eq:deconcatenation}
\Delta[s^{-1}a_1 | \cdots | s^{-1}a_n]
\;=\;
\sum_{i=0}^{n}
[s^{-1}a_1 | \cdots | s^{-1}a_i]
\otimes
[s^{-1}a_{i+1} | \cdots | s^{-1}a_n],
\end{equation}
is a cofree $\Eone$-coalgebra. Here
$\Abar = \ker(\varepsilon \colon \cA \to \omega_X)$ is the
augmentation ideal and $s^{-1}$ is the desuspension, lowering
degree by $1$. The integral kernel is the logarithmic $1$-form
\begin{equation}\label{eq:eta}
\eta_{12} \;=\; d\log(z_1 - z_2)
\;=\; \frac{dz_1 - dz_2}{z_1 - z_2}\,,
\end{equation}
has residue~$1$.  A coordinate change adds an exact logarithmic form
regular along the diagonal, so the residue class is intrinsic. At
three points, Arnold's relation \cite{Arnold69}
\begin{equation}\label{eq:arnold}
\eta_{12} \wedge \eta_{23}
+ \eta_{23} \wedge \eta_{31}
+ \eta_{31} \wedge \eta_{12}
\;=\; 0
\end{equation}
generates all relations in
$H^*(\Conf_n(\mathbf{C}), \mathbf{C})$.  Under the faithful
Fulton--MacPherson jet realization, bar nilpotence and the Borcherds
identity are equivalent: Arnold's relation governs the boundary forms,
and the Borcherds identity governs their OPE coefficients.  This is
the \emph{selection principle}: configuration-space geometry realizes
the chiral identities at every collision stratum.

The binary collision residue of the bar differential is the collision
kernel
\[
r^{\mathrm{coll}}_\cA(z) \;=\;
\Res^{\mathrm{coll}}_{0,2}(\Theta_\cA)
\;\in\; \operatorname{End}_\cA(2)[\![z^{-1}]\!],
\]
where $\Theta_\cA := D_\cA - d_0$ is the Maurer--Cartan element
of the bar differential, satisfying
$d\Theta_\cA + \tfrac{1}{2}[\Theta_\cA, \Theta_\cA] = 0$
automatically because $(D_\cA)^2 = 0$.  A classical $r$-matrix
interpretation requires a Yang--Baxter identity in a specified
representation.  Homotopy coinvariants retain the equivariant
endomorphism data.  A normalized trace functional
\[
 \tau_\cA\colon
 \operatorname{End}_\cA(2)_{h\Sigma_2}\longrightarrow\mathbf C
\]
defines the scalar
\begin{equation}\label{eq:kappa-def}
\kappa(\cA) \;=\; \tau_\cA(r^{\mathrm{coll}}_\cA)
\quad \text{for abelian and scalar families,}
\end{equation}
while for non-abelian affine Kac--Moody
\[
\kappa(V_k(\fg))
\;=\;\tau_{V_k(\fg)}(r^{\mathrm{coll}}_{V_k(\fg)})
      +\dim(\fg)/2.
\]
This is the modular characteristic selected by the trace. Five
structural theorems extract the
$\Sigma_n$-invariant content of the ordered bar. Each theorem
characterizes one facet of $\Theta_\cA$:
\begin{enumerate}[label=\textup{(\Alph*)}]
\item \emph{Reconstruction and Verdier refinement}:
 $R_X(\cA)=\Omega_XB_X(\cA)\xrightarrow{\sim}\cA$ and
 $K_X(\cA)=\mathbb D_{\Ran}B_X(\cA)$ at levels
 $1\leftrightarrow2$.
\item \emph{Quadratic recognition}:
 $q_\cA\colon\cA^{\mathrm i}\xrightarrow{\sim}B_X(\cA)$ under
 $H_{\mathrm{CL}}$.
\item \emph{Stratification}: the $5 \times 5$ $\kappa$-stratification
 matrix on rows $\mathsf{G/L/C/M/B}$ at level~$5$ per stratum.  The
 $K3 \times E$ row has native entries $(0,0,\star,5,24)$, and
 $H_{\mathsf B}$ supplies its Open Heisenberg entry~$3$.
\item \emph{Modular hierarchy}: deformation class, virtual
 $K$-object, signed Hodge shadow, and numerical graph trace, connected
 by the packages $H_D^1$, $H_D^K$, $H_D^{\mathrm{tr}}$, and
 $H_D^{\mathrm{graph}}$.
\item \emph{Family support}: $H_H(\cA;S)$ gives
 $\operatorname{Supp}\ChirHoch^\bullet(\cA)\subseteq S$ at level~$3$.
\end{enumerate}

The chain is typed: A constructs reconstruction and the Verdier
branch; B recognizes the quadratic model; C organizes the five
scalar measurements; D compares deformation, Hodge, and graph
realizations; H computes family-specific cohomological support.

The locus $\mathrm{KSDual} \subset \mathrm{Kosz}(X)$ on which
the Verdier involution $\cA \mapsto \cA^!$ acts by the identity
is the $\mathbb{Z}/2$-fixed sublocus where the four-quadrant
grid $(\text{Open} \cup \text{CY}) \times
(\text{Cat} \cup \text{Chain})$ degenerates into self-dual
form: Theorem~A supplies the Ran unit and counit together with the
separate Verdier branch; Theorem~H supplies the family set~$S$ under
$H_H(\cA;S)$; the
five-archetype label is rigid on the admissible finite-type
curved fixed locus, and
$\kappa_{\mathrm{BKM}}(\Phi_N) = c_N(0)/2$
\cite{Borcherds1998,Gritsenko1999} is the equivariant
signature for $N \in \{1, 2, 3, 4, 6\}$ plus half- and
quarter-integer continuations. Scalar anti-diagonal relations such as
$K=0$ or $c=13$ are duality-orbit sums; the finite-type curved
Verdier--Koszul fixed-point criterion supplies KSDual membership.

\subsection{Prior work}

Five existing programmes supply distinct layers of the problem.
Beilinson--Drinfeld \cite{BD04} construct chiral and factorization
algebras on $\Ran(X)$ together with chiral homology.
Francis--Gaitsgory \cite{FG12} establish enhanced Koszul duality in
the pro-nilpotent Ran category.
Costello--Gwilliam \cite{CG17} develop factorization algebras for
perturbative quantum field theory.
Gui--Li--Zeng \cite{GLZ24} prove a quadratic duality at the
cochain level, a genus-$0$ degree-$2$ result.
Malikov--Schechtman \cite{MS24} construct homotopy chiral algebras
as the appropriate derived input.  The five theorems below organize
the comparison maps among these layers.

\begin{center}
\small
\begin{tabular}{l@{\quad}c@{\quad}c@{\quad}c@{\quad}c@{\quad}c}
\toprule
& \textbf{A} & \textbf{B} & \textbf{C} & \textbf{D} & \textbf{H} \\
\textbf{Programme}
& \footnotesize adjunction
& \footnotesize inversion
& \footnotesize compl.
& \footnotesize genus
& \footnotesize chiral Hochschild \\
\midrule
BD \cite{BD04}
 & $D$ & & & & \\
FG \cite{FG12}
 & $P$ & & & & \\
CG \cite{CG17}
 & & & & & \\
GLZ \cite{GLZ24}
 & & $P_{g=0,\,d=2}$ & & & \\
MS \cite{MS24}
 & & $D$ & & & \\
\textbf{This work}
 & $P+Q$ & $C$ & $C$ & $C$ & $C$ \\
\bottomrule
\end{tabular}
\end{center}

\noindent
{\footnotesize $P$: proved in the cited source; $D$: foundational or
definitional input; $C$: conditional on the displayed comparison
package; $Q$: factorization restriction presently conjectural.  In
column~A, $P+Q$ records the published Ran equivalence together with
its conjectural restriction to the selected factorization
subcategory.}

\subsection{The Heisenberg computation}

The Heisenberg chiral algebra $\cH_k$ is generated by a single
current $J$ of conformal weight $1$ with OPE
$J(z)\,J(w) \sim k/(z-w)^2$.  The normalized coefficient of the
double pole is
\[
 \kappa(\cH_k)=k.
\]
At $k=0$, Theorem~A reconstructs the chart from its ordinary Ran bar.
At $k\ne0$, the second-kind package $H_{\mathrm{curv}}$ supplies the
curved coderived counit.  A quadratic Heisenberg comparison is an
instance of Theorem~B once the finite-window contraction in
$H_{\mathrm{CL}}^{\mathrm{II}}$ is supplied.  Under
the packages of Theorem~D,
\[
 \operatorname{tr}_1\operatorname{Obs}^{\mathrm{def}}_{1,1}(\cH_k)
 =k\lambda_1,
 \qquad
 \operatorname{obs}^{\mathrm{Hdg}}_g(\cH_k)=(-1)^gk\lambda_g,
 \quad g\ge2,
\]
and the scalar-diagonal graph functional has generating series
$k((x/2)/\sin(x/2)-1)$.  For the rank-one even free superboson,
Bakalov--De Sole--Kac compute the bounded vector $(2,1,0,\ldots)$;
a bounded-to-chart quasi-isomorphism transports this vector to
Theorem~H.  Thus the rank-one example fixes the scalar normalization
and displays every comparison map that remains family-dependent.


\subsection{Conventions}

Throughout, $X$ denotes a smooth projective curve over
$\mathbf{C}$, $\cA$ a chiral algebra on $X$ in the sense of
Beilinson--Drinfeld \cite{BD04}.  The stack $\Mbar_{g,n}$ is used in
the stable range $2g-2+n>0$; we write $\Mbar_g=\Mbar_{g,0}$ for
$g\geq2$.
Grading is cohomological: $|d| = +1$. Desuspension lowers degree:
$|s^{-1}v| = |v| - 1$.  On the strictly augmented lane, the bar
complex uses
$\barB(\cA)=T^c(s^{-1}\Abar)$ with
$\Abar=\ker(\varepsilon)$.  On a centrally curved lane, a vacuum
splitting replaces the algebra augmentation.  The second-kind bar
coalgebra carries a degree-two curvature functional~$h_B$; its curved
cobar algebra carries the corresponding element~$m_0$.
The algebraic Fulton--MacPherson compactification
$\overline{\FM}_n(X)$ is the iterated algebraic blowup of $X^n$ along
the diagonal building set \cite{FM94}.  It is smooth, and its boundary
is a simple normal-crossings divisor with components $D_S$ indexed by
subsets $S\subseteq\{1,\ldots,n\}$ of cardinality at least~$2$.
The oriented real blowup of its complex-analytic boundary is the Betti
model with corners.  The algebraic bar uses the first space; its
configuration-space integrals may be realized on the second.


%%% ================================================================
%%% SECTION 2: THE BAR COMPLEX
%%% ================================================================

\section{The bar complex on algebraic curves}
\label{sec:bar-complex}

\subsection{Chiral algebras}

A \emph{chiral algebra} $\cA$ on a smooth curve $X$
\cite{BD04} is a $\cD_X$-module equipped with a chiral bracket
\begin{equation}\label{eq:chiral-bracket}
\mu^{\mathrm{ch}} \colon
j_* j^*(\cA \boxtimes \cA)
\longrightarrow
\Delta_!(\cA),
\end{equation}
where $j \colon X^2 \setminus \Delta \hookrightarrow X^2$ is the
complement of the diagonal and $\Delta_!$ is the $!$-pushforward
along the diagonal embedding. In local coordinates, the chiral
bracket encodes the operator product expansion:
\[
a(z)\,b(w) \;\sim\;
\sum_{n \geq 0} \frac{a_{(n)}b(w)}{(z-w)^{n+1}}.
\]
The zeroth product $a_{(0)}b$ is the Lie bracket (first-order
pole); the first product $a_{(1)}b$ carries the symmetric inner
product (second-order pole). The Borcherds identity
\begin{equation}\label{eq:borcherds}
\sum_{j \geq 0} \binom{m}{j}
(a_{(n+j)}b)_{(m+k-j)}c
\;=\;
\sum_{j \geq 0} (-1)^j \binom{n}{j}
\left(
a_{(m+n-j)}(b_{(k+j)}c)
- (-1)^n\, b_{(n+k-j)}(a_{(m+j)}c)
\right)
\end{equation}
is the defining axiom: it couples all pole orders in a single
constraint.

\subsection{The propagator and the Arnold relation}

The bar differential extracts OPE data from collisions
on $C_2(X) = X \times X \setminus \Delta$.  Poincar\'e residue acts
on logarithmic first-order poles.  The canonical local generator is
the logarithmic derivative
\[
\eta_{12} = d\log(z_1 - z_2)
= \frac{dz_1 - dz_2}{z_1 - z_2}\,,
\qquad \Res_\Delta\, \eta_{12} = 1.
\]
This is a local logarithmic generator.  Its residue is independent of
the chosen coordinate.

Under $z \mapsto w(z)$, the local generator transforms as
$\eta_{12} \mapsto \eta_{12}
+ d\log\frac{w(z_1) - w(z_2)}{z_1 - z_2}$; near the diagonal
the second summand is regular.  Thus the class of $\eta_{12}$ in the
logarithmic normal quotient, and hence its Poincar\'e residue, glues
canonically; a global propagator requires the additional choice of a
connection or prime-form normalization.

At three points, Arnold's relation \cite{Arnold69}
\[
\eta_{12} \wedge \eta_{23}
+ \eta_{23} \wedge \eta_{31}
+ \eta_{31} \wedge \eta_{12}
= 0
\]
restates the partial-fractions identity
\[
\frac{1}{(z_1 - z_2)(z_2 - z_3)}
+ \frac{1}{(z_2 - z_3)(z_3 - z_1)}
+ \frac{1}{(z_3 - z_1)(z_1 - z_2)}
= 0.
\]
Arnold proved that, replicated across all triples
$\{i,j,k\} \subset \{1, \ldots, n\}$, this generates the
complete ideal of relations in
$H^*(\Conf_n(\mathbf{C}), \mathbf{C})$.

\subsection{The ordered bar complex}

\begin{definition}[Fulton--MacPherson realization of the ordered bar;
\ClaimStatusConditional]\label{def:ordered-bar}
Let $\cA$ be a strictly augmented chiral algebra on a smooth curve
$X$, let $\Abar=\ker(\varepsilon)$, and let
$j_n\colon\Conf_n(X)\hookrightarrow X^n$ and
$\pi_n\colon\overline{\FM}_n(X)\to X^n$ be the open inclusion and
blowdown map.  The enhanced bar $B_X(\cA)$ is formed first in
$D(\Ran X)$.  The Fulton--MacPherson jet-realization package
$H_{\mathrm{FM}}^{\mathrm{jet}}$ identifies its arity-$n$ chain model
with
\begin{equation}\label{eq:bar-full}
 B^{\mathrm{ord}}_{X,n}(\cA)
 \simeq
 R\Gamma\!\left(
  \overline{\FM}_n(X),
  \operatorname{DR}_{\log}\!\left(
   \pi_n^!j_{n,*}j_n^*(s^{-1}\Abar)^{\boxtimes n}
  \right)
 \right),
\end{equation}
with the full logarithmic de Rham complex and its total grading.  The
summand with $\Omega^{n-1}_{\log}$ is one graded piece of this complex.
The ordered object precedes passage to homotopy
$\Sigma_n$-coinvariants.  Here $s^{-1}$ is desuspension and
$|d|=+1$.
\end{definition}

The bar differential decomposes into three pieces:
\begin{equation}\label{eq:bar-diff}
d_{\barB} = d_{\mathrm{int}} + d_{\mathrm{res}} + d_{\mathrm{dR}}.
\end{equation}
The internal differential $d_{\mathrm{int}}$ acts on each tensor
factor of $\cA$. The de Rham differential $d_{\mathrm{dR}}$ acts
on logarithmic forms on $\overline{\FM}_n(X)$.  The collision
differential combines two pieces of geometry.  The Poincar\'e residue
of $\eta_{ij}$ records the boundary incidence, while the normal-jet
filtration of $j_*j^*(\cA\boxtimes\cA)$ records the OPE pole order.
On the $q$-th principal-part graded piece, $q\geq1$, their contraction
is
\begin{equation}\label{eq:res-diff}
 \operatorname{gr}^{q}_{D_{ij}}d_{\mathrm{res}}^{(ij)}
 \bigl(a_1\otimes\cdots\otimes a_n\otimes
       (z_i-z_j)^{-q}\otimes\eta_{ij}\wedge\omega'\bigr)
 =\pm\,a_{i(q-1)}a_j\otimes
       \widehat{a_1\otimes\cdots\otimes a_n}^{\,ij}
       \otimes\omega'.
\end{equation}
Here the hat with superscript $ij$ retains every tensor factor except
the $i$-th and $j$-th.  In particular, a pole of order~$q$ produces
the mode $a_{(q-1)}b$.  The jet bundle makes this associated-graded
description coordinate-covariant.

\subsection{Why $d^2 = 0$: the Borcherds selection principle}

\begin{theorem}[Bar nilpotence; \ClaimStatusConditional]
\label{thm:bar-nilpotence}
Under $H_{\mathrm{FM}}^{\mathrm{jet}}$, the Borcherds identity
\eqref{eq:borcherds} implies $d_{\barB}^2=0$.  If the jet realization
is faithful on the normal-principal-part filtration, the converse also
holds.
\end{theorem}

\begin{proof}
Expand $(d_{\mathrm{int}} + d_{\mathrm{res}} + d_{\mathrm{dR}})^2$
into nine terms. Three vanish individually:
$d_{\mathrm{int}}^2 = 0$ (cochain complex),
$d_{\mathrm{dR}}^2 = 0$ (de Rham), and
$\{d_{\mathrm{int}}, d_{\mathrm{dR}}\} = 0$ (Leibniz). The
remaining terms are controlled by index overlap in the residue
operator:

\emph{Disjoint indices}
($\{i,j\} \cap \{k,\ell\} = \emptyset$): residue operators
commute; cancellation is by the Koszul sign rule.  The OPE input
enters at overlapping collision strata.

\emph{One shared index} (say $j = k$): three boundary strata meet
at the codimension-$2$ locus $D_{ij\ell}$ of triple collision.
The Borcherds identity makes the cyclic sum of iterated residues
vanish.  Arnold's relation contracts the differential-form factor,
while the Borcherds identity contracts the OPE factor.  Faithfulness
of the normal-jet realization gives the converse.

\emph{Same pair} ($\{i,j\} = \{k,\ell\}$):
$(\Res_{D_{ij}})^2 = 0$ because the first Poincar\'e residue removes
the unique logarithmic normal factor along \(D_{ij}\).
\end{proof}

\begin{remark}[Selection table]\label{rem:selection-table}
The bar construction asks the same question at every level
of geometry:
\begin{center}
\small
\begin{tabular}{lll}
\textbf{Base} & \textbf{Condition for $d^2 = 0$}
 & \textbf{Selected objects} \\
\hline
Point & Jacobi identity & Lie algebras \\
Curve $X$ & Borcherds identity & Chiral algebras \\
$X$ over $\Mbar_{g,n}$ in the stable range & $+$ period corrections
 & Modular Koszul algebras
\end{tabular}
\end{center}
Each row is forced by the geometry of the next.
\end{remark}

\subsection{Two coproducts}

The ordered bar carries the \emph{deconcatenation
coproduct}~\eqref{eq:deconcatenation}: $n+1$ terms, coassociative and
$\Eone$, reflecting the linear ordering. The homotopy-$\Sigma_n$
descent is
\[
 B^\Sigma_{X,n}(\cA)
 :=\bigl(B^{\mathrm{ord}}_{X,n}(\cA)\bigr)_{h\Sigma_n}.
\]
Under the symmetric descent package it inherits the coshuffle
coproduct; on the tensor associated graded this coproduct has $2^n$
summands and is cocommutative. The
ordered bar is the level-$2$ object before symmetrisation; the symmetric
bar is the completed coinvariant descent. In the braided case this is
the $L_R$-twisted descent, and the formal conductor fibre records the
non-coinvariant $R$-matrix data.

\subsection{From genus $0$ to all genera}

At genus $g \geq 1$, the prime form globalizes the local difference.
A normalized logarithmic kernel is
$d\log E(z_1, z_2)$, where $E(z,w)$ is the prime form on a
Riemann surface of genus $g$ (a section of
$K^{-1/2} \boxtimes K^{-1/2}$, weight $1$).  The second-kind bar
coalgebra carries a curvature functional, and its curved cobar
presentation carries an element $m_0^{(g)}$ whose scalar realization
is
\begin{equation}\label{eq:curvature}
\operatorname{tr}_{\mathrm{diag}}\!\bigl(m_0^{(g)}\bigr)
\;=\; \kappa(\cA) \cdot \omega_g,
\end{equation}
after scalar diagonal projection of the matrix-valued ordered
curvature residue
$r^{\mathrm{coll}}_\cA(z)\cdot\omega_g\in
\operatorname{End}_\cA(2)\otimes H^{1,1}(\cC_g)$. The raw chain
identity is $m_1^2=[m_0^{(g)},-]$.

Under the modular Maurer--Cartan package, the \emph{total} bar
differential
$D_\cA = \sum_{g \geq 0} \hbar^g\, d_\cA^{(g)}$ incorporates the
curvature and satisfies $D_\cA^2 = 0$ in the completed ambient.

\begin{definition}[Modular characteristic]\label{def:kappa}
The \emph{modular characteristic} \(\kappa(\cA)\) is the normalized
degree-two collision trace. In scalar or abelian cases it equals the
value of \(\tau_\cA\) on the collision kernel. For nonabelian affine
Kac--Moody algebras in the
trace-form normalization,
\[
\kappa(V_k(\fg))
\;=\;
\tau_{V_k(\fg)}(r^{\mathrm{coll}})+\frac{\dim(\fg)}{2}
\;=\;
\frac{\dim(\fg)(k+h^\vee)}{2h^\vee},
\]
where the second summand is the Sugawara simple-pole contribution.
\end{definition}

\begin{definition}[Maurer--Cartan element]\label{def:mc}
The \emph{Maurer--Cartan element} is
$\Theta_\cA := D_\cA - d_0 \in \gAmod$, the positive-genus
part of the bar differential. Since $D_\cA^2 = 0$, it satisfies
\begin{equation}\label{eq:mc}
d_0\, \Theta_\cA + \tfrac{1}{2}[\Theta_\cA, \Theta_\cA] = 0
\end{equation}
in the completed modular ambient.  The five theorems characterize
its reconstruction, duality, trace, and support realizations.
\end{definition}


%%% ================================================================
%%% SECTION 3: THEOREM A
%%% ================================================================

\section{Theorem A: Ran reconstruction and Verdier refinement}
\label{sec:theorem-A}

\subsection{Statement}

The full bar object carries three canonical exits, each with its own
target and hypothesis package:
\begin{equation}\label{eq:three-functors}
\begin{array}{lcll}
R_X(\cA):=\Omega_XB_X(\cA)&\xrightarrow{\ \varepsilon_\cA\ }&\cA
 &\text{(Ran reconstruction),}\\[3pt]
\cA^{\mathrm i}&\xrightarrow{\ q_\cA\ }&B_X(\cA)
 &\text{(quadratic compression),}\\[3pt]
K_X(\cA):=\mathbb D_{\Ran}B_X(\cA)
&&
 &\text{(Verdier--Koszul algebra).}
\end{array}
\end{equation}
Reconstruction returns the source algebra.  Quadratic compression
compares a presentation-dependent coalgebra with the full bar.
Verdier duality transports the full bar coalgebra to an algebra.

\begin{definition}[The Ran ambient and four hypothesis packages]
\label{def:A-bang-infty}
Let
$\mathcal C_X=(D(\Ran X),\otimes^{\mathrm{ch}})$.
The proof of Francis--Gaitsgory's Theorem~5.1.1 makes this chiral
monoidal category pro-nilpotent, and their Proposition~4.1.2 applies
to the reduced associative operad.  The package
$H_{\mathrm{fact}}$ asserts closure of
the chosen associative factorization subcategories under enhanced
bar, enhanced cobar, collision base change, units, and completed
inverse limits.  The package
$H_{\mathrm{CL}}(\cA,\cA^{\mathrm i},\tau_{\mathrm i})$ supplies a
connected positive-weight quadratic coalgebra $\cA^{\mathrm i}$,
its canonical twisting morphism
$\tau_{\mathrm i}\colon\cA^{\mathrm i}\to\cA$, finite-window
comparison, and acyclicity of the associated twisted tensor products.
The package $H_{\mathrm{conv}}$ realizes
the completed chain model through exhaustive complete filtrations
and cohomologically Mittag--Leffler inverse systems.  The package
$H_{\mathrm{VD}}$ places
$B_X(\cA)$ in a constructible or holonomic dualizable subcategory and
supplies the Kunneth and base-change maps used by Verdier duality.
\end{definition}

\begin{theorem}[Theorem A: typed bar--cobar reconstruction;
\ClaimStatusProvedElsewhere]
\label{thm:bar-cobar-adjunction}
\textsc{Type signature.} Quadrant: Open. Presentation: chain-level
$\infty$-category. Level: $1\leftrightarrow2$.  Ambient:
the pro-nilpotent Francis--Gaitsgory chiral Ran category with the
reduced associative operad.

For the reduced associative operad in
$\mathcal C_X=(D(\Ran X),\otimes^{\mathrm{ch}})$, enhanced bar and
cobar are mutually inverse equivalences
\cite[Proposition~4.1.2, together with the pro-nilpotence argument in
the proof of Theorem~5.1.1]{FG12}.
For an augmented algebra $\cA$ and a divided-power conilpotent
coalgebra $C$, the unit and counit are equivalences
\[
 C\xrightarrow{\ \sim\ }B_X\Omega_X(C),
 \qquad
 \Omega_XB_X(\cA)\xrightarrow{\ \sim\ }\cA.
\]
\end{theorem}

\begin{conjecture}[Factorization restriction; \ClaimStatusConjectured]
\label{conj:standalone-factorization-restriction}
Under $H_{\mathrm{fact}}$, enhanced bar and cobar preserve the
factorization equivalences and collision maps, and the ambient
Francis--Gaitsgory equivalence restricts to the associative
factorization subcategories used in this volume.
\end{conjecture}

\begin{proposition}[Quadratic compression; \ClaimStatusConditional]
\label{prop:standalone-quadratic-compression}
Under $H_{\mathrm{CL}}(\cA,\cA^{\mathrm i},\tau_{\mathrm i})$, the canonical map
\[
 q_\cA\colon\cA^{\mathrm i}\xrightarrow{\ \sim\ }B_X(\cA)
\]
is a quasi-isomorphism.  Local finite-dimensionality forms the
quadratic dual algebra $\cA^!=(\cA^{\mathrm i})^\vee$.
\end{proposition}

\begin{proposition}[Verdier refinement; \ClaimStatusConditional]
\label{prop:standalone-verdier-refinement}
Under $H_{\mathrm{VD}}$, Verdier duality transports the coalgebra
$B_X(\cA)$ to the algebra
\[
 K_X(\cA):=\mathbb D_{\Ran}B_X(\cA),
\]
and biduality gives
$B_X(\cA)\xrightarrow{\sim}\mathbb D_{\Ran}K_X(\cA)$.
Together with quadratic compression, this identifies $K_X(\cA)$
with the continuous Verdier realization of $\cA^!$.
\end{proposition}

\begin{definition}[The KSDual fixed locus; \ClaimStatusDefinitional]
\label{def:standalone-ksdual-fixed-locus}
A point belongs to $\mathrm{KSDual}$ when it carries a chosen
equivalence $\sigma_\cA\colon\cA\xrightarrow{\sim}K_X(\cA)$
compatible with the evaluation and coevaluation maps.  This is the
$\mathbb Z/2$-fixed sublocus of the specified duality correspondence.
\end{definition}

\begin{remark}[Reconstruction, compression, and duality]
\label{rem:T1-threeway}
The counit $R_X(\cA)\to\cA$ is the reconstruction morphism.  The map
$q_\cA$ compares the quadratic coalgebra with the full bar.  The
algebra $K_X(\cA)$ is the Verdier image of the full bar.  A comparison
$R_X(\cA)\simeq K_X(\cA)$ is precisely reconstruction composed with
a self-duality equivalence
$\sigma_\cA\colon\cA\xrightarrow{\sim}K_X(\cA)$.
\end{remark}

\begin{remark}[Completed realization for unbounded weights]
\label{rem:T1a-ambient-classM}
An unbounded conformal-weight grading is represented by its inverse
system of finite-weight windows.  The package $H_{\mathrm{conv}}$
requires cohomological Mittag--Leffler transition maps and the
comparison with the derived inverse limit.  This completed realization
is the ambient used for the Virasoro and principal $\cW$ charts.
\end{remark}

\subsection{Proof}

Francis--Gaitsgory prove enhanced operadic bar--cobar equivalence in
the pro-nilpotent Ran category.  The maps in $H_{\mathrm{fact}}$
carry that equivalence to the chosen factorization subcategories.
The twisting-morphism fundamental theorem identifies acyclicity of
$\cA\otimes_{\tau}\cA^{\mathrm i}$ with the quasi-isomorphism
$q_\cA$.  Finally, $H_{\mathrm{VD}}$ transports the coalgebra
structure through Verdier duality and supplies biduality.

\begin{example}[Square-zero local model]\label{ex:thmA-heis}
Let $A=k[\epsilon]/(\epsilon^2)$ with its standard augmentation.  The
reduced bar has one basis word $(s^{-1}\epsilon)^{\otimes r}$ in every
positive length and zero bar differential.  Its deconcatenation
coalgebra is the quadratic coalgebra $A^{\mathrm i}$, so $q_A$ is an
isomorphism.  The cobar counit $\Omega B(A)\to A$ is the standard
free resolution.  Finite-dimensional duality on each length gives
the completed quadratic-dual algebra.  This example exhibits
reconstruction, quadratic compression, and duality as three distinct
maps.
\end{example}


%%% ================================================================
%%% SECTION 4: THEOREM B
%%% ================================================================

\section{Theorem B: quadratic Koszul recognition}
\label{sec:theorem-B}

\subsection{Statement}

Theorem~A reconstructs an algebra from its full bar in the Ran
ambient.  Theorem~B characterizes the quadratic coalgebras that
represent this full bar object.

\begin{theorem}[Theorem B: quadratic comparison criterion;
\ClaimStatusConditional]
\label{thm:bar-cobar-inversion}
\textsc{Type signature.} Quadrant: Open. Presentation: chain-level
in the associative factorization category. Level: $1\leftrightarrow2$.
Hypothesis package: $H_{\mathrm{CL}}(\cA,\cA^{\mathrm i},\tau_{\mathrm i})$.

Let $\cA=T_X(V)/(R)$ be a connected positive-weight quadratic
presentation and let
\[
 \cA^{\mathrm i}=C_X(s^{-1}V,s^{-2}R),
 \qquad
 \tau_{\mathrm i}\colon\cA^{\mathrm i}\longrightarrow\cA
\]
be its quadratic coalgebra and canonical twisting morphism.  The
package $H_{\mathrm{CL}}$ supplies factorization descent, compatible
finite-weight windows, convergence of the bar filtration, and the
twisted tensor products in the selected completed ambient.  Under
this package, the following conditions are equivalent:
\begin{enumerate}[label=\textup{(B\arabic*)}]
\item the canonical coalgebra map
$q_\cA\colon\cA^{\mathrm i}\to B_X(\cA)$ is a quasi-isomorphism;
\item the cobar comparison
$\Omega_X(\cA^{\mathrm i})\to\cA$ is a quasi-isomorphism;
\item the left and right twisted tensor products associated with
$\tau_{\mathrm i}$ are acyclic in positive weight.
\end{enumerate}
The additional package $H_{\mathrm{PBW}}^{\mathrm{det}}$ consists of
strict filtration-preserving comparison maps, an associated-graded
quadratic identification, a splitting that detects cohomology, and a
specified $E_2$ page.  Under
$H_{\mathrm{CL}}\cup H_{\mathrm{PBW}}^{\mathrm{det}}$, the preceding
conditions are also equivalent to
\begin{enumerate}[label=\textup{(B4)}]
\item the PBW/Li--bar spectral sequence collapses at $E_2$, and
$E_\infty$ is supported on the Koszul diagonal.
\end{enumerate}
Whenever these conditions hold, local finite-dimensionality forms
the quadratic dual algebra
$\cA^!=(\cA^{\mathrm i})^\vee$, and the Verdier package of
Theorem~A compares it with $K_X(\cA)$.
\end{theorem}

\subsection{Proof architecture}

The twisting-morphism fundamental theorem identifies the acyclicity
of the two twisted tensor products with the quasi-isomorphisms in
conditions~\textup{(B1)} and~\textup{(B2)}.  The strict filtration
and detecting splitting in $H_{\mathrm{PBW}}^{\mathrm{det}}$ identify
the associated graded cone of $q_\cA$ with the off-diagonal part of
the quadratic Koszul complex.  Collapse at $E_2$ and diagonal support
identify the cone with an acyclic complex.  The same argument in each finite-weight
window, followed by the cohomological Mittag--Leffler comparison,
gives the completed statement.

\begin{example}[Polynomial algebra]\label{ex:thmB-heis}
For $A=k[x]$ with $x$ in positive weight, the quadratic coalgebra is
the exterior coalgebra on $s^{-1}x$.  The map
$A^{\mathrm i}\to B(A)$ includes the standard Koszul cycles, and the
Koszul contraction gives a quasi-isomorphism.  Equivalently,
$A\otimes_\tau A^{\mathrm i}$ is the two-term Koszul resolution of
the augmentation module.  This is the first non-trivial instance of
all four conditions.
\end{example}

\begin{remark}[Recognition boundary]\label{rem:off-koszul}
For a family whose cone $\operatorname{Cone}(q_\cA)$ carries
cohomology, the quadratic coalgebra records a proper compression of
the full bar.  The full Ran counit of Theorem~A and the quadratic
comparison of Theorem~B retain their separate domains.  For an
unbounded conformal-weight family, a comparison between direct-sum
and completed product totalizations additionally requires
finite-window contractions compatible with the transition maps.
\end{remark}


%%% ================================================================
%%% SECTION 5: THEOREM C
%%% ================================================================

\section{Theorem C: the \(5\times5\) stratification atlas}
\label{sec:theorem-C}

\subsection{Five measurements}

The symbol \(\kappa\) occurs in five native categories.  We therefore
write
\[
 \boldsymbol\kappa(\cA)
 :=
 \left(
 \kappa_{\mathrm{cat}}(\cA),
 \kappa_{\mathrm{ch}}^{\mathrm{Hodge}}(\cA),
 \kappa_{\mathrm{ch}}^{\mathrm{Heis}}(\cA),
 \kappa_{\mathrm{BKM}}(\cA),
 \kappa_{\mathrm{fiber}}(\cA)
 \right).
\]
These coordinates are respectively categorical, Hodge-graded
chiral, Heisenberg-normalized chiral, Borcherds, and topological.
An equality between two coordinates is represented by a comparison
map between their source objects.

\subsection{Three comparison tiers}

Theorem~C has three typed tiers.

\begin{enumerate}[label=\textsc{C\arabic*.}]
\item The strict flat fibre--centre comparison is the map
\[
 Z\!\left(\Omega_XB_X(\cA)\right)
 \xrightarrow{\ Z(\varepsilon_\cA)\ }
 Z(\cA)=H^0C^\bullet_{\mathrm{ch}}(\cA,\cA)
\]
on a chirally Koszul datum.

\item A represented Verdier involution, perfectness, and a
nondegenerate anti-invariant pairing give the homotopy eigenspace
decomposition and its scalar trace shadow.

\item A shifted-symplectic/BV package upgrades the represented
eigenspaces to a Lagrangian intersection in the global deformation
space.  Its data include the BV Laplacian, the QME anomaly class, and
a chosen null-homotopy of that class whenever QME anomaly cancellation
is asserted.
\end{enumerate}

The scalar relation belongs to the trace realization of \textsc{C1}
together with Theorem~D.  The shifted-symplectic assertion belongs to
\textsc{C2} and carries its own construction.

\begin{definition}[Family comparison package for Theorem C;
\ClaimStatusDefinitional]
\label{def:standalone-theorem-C-package}
For a family \(\mathcal F\), the package \(H_C(\mathcal F)\) consists
of:

\begin{enumerate}[label=\textup{(C\arabic*)}]
\item a source object and trace functional for every displayed
coordinate of \(\boldsymbol\kappa\);
\item a represented Verdier involution on the selected perfect
ambient and its evaluation--coevaluation maps;
\item an anti-invariant pairing whenever a complementarity
decomposition is asserted;
\item a BV and shifted-symplectic package whenever the C2 enhancement
is asserted;
\item cross-quadrant realization maps for every Open--CY equality.
\end{enumerate}
\end{definition}

\subsection{Statement}

\begin{theorem}[Theorem C: census-scoped \(5\times5\) stratification;
\ClaimStatusConditional]
\label{thm:complementarity}
\textsc{Type signature.} Quadrants: Open, with a conditional
Open--CY row. Presentations: categorical, chain, scalar trace, and
topological. Levels: \(0\)--\(5\). Hypothesis package:
\(H_C(\mathcal F)\).

For every family carrying \(H_C(\mathcal F)\), the five measurements
form the row \(\boldsymbol\kappa(\mathcal F)\).  The finite atlas uses
the family labels
\[
 \mathsf G,\qquad\mathsf L,\qquad\mathsf C,\qquad
 \mathsf M,\qquad\mathsf B.
\]
The first four label the selected Vol~I standard-family census.  The
fifth is the conditional \(K3\times E\) comparison row.  Under its
realization package \(H_{\mathsf B}\), that row is
\[
 \boldsymbol\kappa(\mathsf B)=(0,0,3,5,24).
\]
Here \(\mathrm{wt}(\Delta_5)=5\) \cite{Gritsenko1999} and
\(e_{\mathrm{top}}(K3)=24\) \cite{Mukai1987} are exact native
calculations.  In the stated normalizations the categorical and Hodge
entries are the native zeros.  The Open-quadrant Heisenberg entry~$3$
is transported by the comparison
\(\eta_{K3,E}\) in \(H_{\mathsf B}\).  The holomorphic Euler
characteristic \(\chi(\mathcal O_{K3})=2\) remains a distinct
categorical invariant.

Suppose the represented Verdier package gives a perfect involution
\(\iota_{\mathrm{VD}}\) on a genus-\(g\) ambient complex
\(\mathbf C_g(\cA)\).  Then its idempotents
\[
 e_\pm=\frac12(\id\pm\iota_{\mathrm{VD}})
\]
give the homotopy decomposition
\begin{equation}\label{eq:lagrangian}
 \mathbf C_g(\cA)\simeq
 \operatorname{im}(e_+)\oplus\operatorname{im}(e_-).
\end{equation}
When the anti-invariant perfect pairing identifies these two images
with the chart and dual-chart sectors, cohomology gives
\begin{equation}\label{eq:complementarity-cohom}
 H^\bullet(\mathbf C_g(\cA))
 \cong Q_g(\cA)\oplus Q_g(\cA^!).
\end{equation}
A scalar relation
\begin{equation}\label{eq:standalone-scalar-comparison}
 \kappa(\cA)+\kappa(\cA^!)=K_{\mathcal F}
\end{equation}
is the trace of a represented family comparison.  Its constant
\(K_{\mathcal F}\) belongs to that family datum.

Finally, the C2 package upgrades
\eqref{eq:lagrangian} to the corresponding shifted-symplectic
Lagrangian statement.
\end{theorem}

\begin{proof}
The five coordinates are the trace functionals in
\(H_C(\mathcal F)\).  The idempotent identities
\(e_\pm^2=e_\pm\), \(e_+e_-=0\), and \(e_++e_-=\id\)
give \eqref{eq:lagrangian}.  Perfectness and the represented pairing
identify the two images and give
\eqref{eq:complementarity-cohom}.  Applying the scalar trace to the
represented Verdier comparison gives
\eqref{eq:standalone-scalar-comparison}.  The final clause is the
shifted-symplectic structure supplied by the C2 package.
\end{proof}

\begin{proposition}[Worked scalar pairings;
\ClaimStatusConditional]
\label{prop:complementarity-sums}
A represented trace comparison gives the familiar scalar formulas
\[
 \kappa(\beta\gamma_\lambda)
 +\kappa(bc_\lambda)=0,
 \qquad
 \kappa(\operatorname{Vir}_c)
 +\kappa(\operatorname{Vir}_{26-c})=13.
\]
The first follows from
\(\kappa(\beta\gamma_\lambda)=6\lambda^2-6\lambda+1\) and
\(\kappa(bc_\lambda)=-\kappa(\beta\gamma_\lambda)\).
The second follows from
\(\kappa(\operatorname{Vir}_c)=c/2\).
At \(c=13\), the two Virasoro scalar coordinates agree.
\end{proposition}

\begin{remark}[KSDual fixed locus]
A KSDual point carries a chosen equivalence
\[
 \sigma_\cA\colon\cA\xrightarrow{\sim}K_X(\cA)
\]
compatible with the represented pairing.  Its scalar coordinate is
then fixed by the induced trace comparison.  The object-level
equivalence and its scalar trace occupy successive levels of the
tower.
\end{remark}

%%% ================================================================
%%% SECTION 6: THEOREM D
%%% ================================================================

\section{Theorem D: the modular characteristic hierarchy}
\label{sec:theorem-D}

\subsection{The four typed objects}

The genus expansion carries four successive realizations of the
modular anomaly.  Their ambient groups determine the meaning of every
formula.

\begin{definition}[Deformation class, virtual object, Hodge shadow,
and graph trace; \ClaimStatusDefinitional]
\label{def:standalone-theorem-D-objects}
For $g\geq2$, let
\[
 \mathrm{Def}_g(\cA)
 :=R\Gamma\!\left(
 \Mbar_g,\mathcal E nd(\bar B^{(g)}_{\mathrm{flat}}(\cA))
 \right).
\]
The native Maurer--Cartan obstruction is
\[
 \operatorname{Obs}^{\mathrm{def}}_g(\cA)
 \in H^2(\mathrm{Def}_g(\cA)).
\]
A perfect scalar realization gives a virtual object
\[
 \mathfrak O_g^K(\cA)\in K^0(\Mbar_g)\otimes\mathbb C,
\]
whose codimension-\(g\) Hodge shadow is
\[
 \operatorname{obs}^{\mathrm{Hdg}}_g(\cA)
 :=\operatorname{ch}_g(\mathfrak O_g^K(\cA))
 \in H^{2g}(\Mbar_g,\mathbb C).
\]
A stable-graph functional gives the numerical coefficient
\(F_g(\cA)\in\mathbb C\) from the marked stable stack
\(\Mbar_{g,1}\), for $g\geq1$.  At genus one, the pointed deformation
complex
\[
 \mathrm{Def}_{1,1}(\cA)
 :=R\Gamma\!\left(
 \Mbar_{1,1},\mathcal E nd(\bar B^{(1,1)}_{\mathrm{flat}}(\cA))
 \right)
\]
is the ambient of
\(\operatorname{Obs}^{\mathrm{def}}_{1,1}(\cA)
\in H^2(\mathrm{Def}_{1,1}(\cA))\), the source of the normalized
one-loop trace.
\end{definition}

\begin{lemma}[Higher-order Maurer--Cartan obstruction class;
\ClaimStatusProvedHere]
\label{lem:standalone-MC-obstruction}
Let \(L\) be a complete filtered dg Lie algebra, let
\(\theta_0\in L^1\) be a Maurer--Cartan element, and fix \(g\geq2\).
Suppose
\(\theta_{<g}(t)=\theta_0+\sum_{r=1}^{g-1}t^r\theta_r\)
satisfies the Maurer--Cartan equation modulo \(t^g\).  Then
\[
 o_g:=\frac12
 \sum_{\substack{i+j=g\\1\leq i,j\leq g-1}}
 [\theta_i,\theta_j]
\]
is a cocycle for \(d_{\theta_0}=d+[\theta_0,-]\).  A coefficient
\(\theta_g\) extends the solution through order \(t^g\) precisely when
\([o_g]=0\) in \(H^2(L,d_{\theta_0})\).
\end{lemma}

\begin{proof}
Write \(F(\theta)=d\theta+\tfrac12[\theta,\theta]\).  The coefficients
of \(F(\theta_{<g})\) below order~$g$ vanish.  The coefficient of
$t^g$ in the Bianchi identity \(d_\theta F(\theta)=0\) is therefore
\(d_{\theta_0}o_g=0\).  Adding $t^g\theta_g$ changes this coefficient
to \(o_g+d_{\theta_0}\theta_g\), which gives the stated equivalence.
\end{proof}

\begin{remark}[The genus-one anchor]
The sum defining \(o_g\) begins at \(g=2\).  The genus-one class in
Theorem~D arises from the normalized one-loop collision trace; the
datum \(H_D^1\) identifies that trace with
\(\kappa(\cA)\lambda_1\).
\end{remark}

\subsection{Statement}

\begin{theorem}[Theorem D: modular characteristic hierarchy;
\ClaimStatusConditional]
\label{thm:genus-universality}
\textsc{Type signature.} Quadrant: Open. Presentation: modular
deformation complex, perfect \(K\)-theory realization, and stable
graphs. Levels: \(4\to5\). Hypothesis package:
\[
 H_D=(H_D^1,H_D^K,H_D^{\mathrm{tr}},H_D^{\mathrm{graph}}).
\]

Let \(\cA=A_b\) be a cyclic modular Koszul chart algebra and let
\(\kappa(\cA)\) be its normalized degree-two collision trace.

\begin{enumerate}[label=\textup{(D\arabic*)}]
\item Under \(H_D^1(\cA)\), the normalized genus-one trace is
\begin{equation}\label{eq:theorem-D-genus-one-standalone}
 \operatorname{tr}_1
 \operatorname{Obs}^{\mathrm{def}}_{1,1}(\cA)
 =\kappa(\cA)\lambda_1
 \in H^2(\Mbar_{1,1}).
\end{equation}

\item For $g\geq2$, under \(H_D^K(\cA,g)\), the perfect virtual
object on \(\Mbar_g\) is
\[
 \mathfrak O_g^K(\cA)
 =\kappa(\cA)\lambda_{-1}(\mathbb E_g),
 \qquad
 \lambda_{-1}(\mathbb E_g)
 =\sum_{i=0}^g(-1)^i[\Lambda^i\mathbb E_g].
\]
Its top Hodge projection is
\begin{equation}\label{eq:theorem-D-hodge-standalone}
 \operatorname{obs}^{\mathrm{Hdg}}_g(\cA)
 =(-1)^g\kappa(\cA)\lambda_g.
\end{equation}
The package \(H_D^{\mathrm{tr}}(\cA,g)\) supplies the comparison
from the scalar realization of
\(\operatorname{Obs}^{\mathrm{def}}_g(\cA)\) to
\(\mathfrak O_g^K(\cA)\).

On \(\Mbar_{1,1}\), the analogous virtual object has degree-one
character \(-\kappa(\cA)\lambda_1\).  This Hodge character and the
positively normalized one-loop trace in~\textup{(D1)} are two
realizations connected only by their stated comparison data.

\item For $g\geq1$, under \(H_D^{\mathrm{graph}}(\cA,g)\), the
stable-graph
functional is
\begin{equation}\label{eq:theorem-D-graph-standalone}
 F_g(\cA)
 =\kappa(\cA)\lambda_g^{\mathrm{FP}}
 +\delta F_g^{\mathrm{cross}}(\cA),
 \qquad
 \lambda_g^{\mathrm{FP}}
 =\int_{\Mbar_{g,1}}\psi_1^{2g-2}\lambda_g.
\end{equation}
The scalar-diagonal graph package gives
\(\delta F_g^{\mathrm{cross}}(\cA)=0\).
For a multi-weight chart, the second term collects stable graphs whose
edges meet distinct conformal-weight channels.
\end{enumerate}
\end{theorem}

\begin{remark}[Scope and normalizations]
\label{rem:thmD-scope}
\label{rem:genus-1-anchor}
The genus-one trace in
\eqref{eq:theorem-D-genus-one-standalone} has positive sign by its
loop normalization.  The Hodge shadow in
\eqref{eq:theorem-D-hodge-standalone} carries the sign of the
alternating exterior power.  The graph functional in
\eqref{eq:theorem-D-graph-standalone} is numerical.  The comparison
packages connect these three realizations while preserving their
ambient degrees.
\end{remark}

\subsection{The Chern-root calculation}
\label{subsec:thmD-proof}

\begin{proposition}[Top Hodge projection;
\ClaimStatusProvedHere]
\label{prop:scalar-obs-hodge-euler}
\label{lem:T4-shared-input}
Let \(E\) be a rank-\(g\) vector bundle.  Then
\[
 \operatorname{ch}_g(\lambda_{-1}(E))
 =(-1)^g c_g(E).
\]
Consequently \(H_D^K(\cA,g)\) gives
\(\operatorname{obs}^{\mathrm{Hdg}}_g(\cA)
=(-1)^g\kappa(\cA)\lambda_g\).
\end{proposition}

\begin{proof}
On a splitting space, let \(x_1,\ldots,x_g\) be the Chern roots of
\(E\).  Then
\[
 \operatorname{ch}(\lambda_{-1}(E))
 =\prod_{i=1}^g(1-e^{x_i})
 =(-1)^g x_1\cdots x_g
  +\text{terms of codimension greater than }g.
\]
Since \(x_1\cdots x_g=c_g(E)\), the codimension-\(g\) component is the
displayed class.
\end{proof}

\begin{remark}[The comparison problem]
\label{rem:T4-shared-input}
\label{rem:clutching-gap-resolved}
A construction of \(H_D^K\) consists of a perfect scalar virtual bar
object and an equality
\[
 \beta_g(\cA)\colon
 \mathfrak O_g^K(\cA)
 \xrightarrow{\ \sim\ }
 \kappa(\cA)\lambda_{-1}(\mathbb E_g)
\]
compatible with clutching and modular trace.  Grothendieck--Riemann--
Roch computes the characteristic class after this comparison is
given.  Clutching tests its boundary compatibility.  The map
\(\beta_g(\cA)\) is the load-bearing higher-genus construction.
\end{remark}

\subsection{The Faber--Pandharipande scalar}

Faber--Pandharipande's $\lambda_g$ formula
\cite[Section~0.2]{FP03} gives
\[
 \lambda_g^{\mathrm{FP}}
 =\frac{2^{2g-1}-1}{2^{2g-1}}\,
  \frac{|B_{2g}|}{(2g)!}.
\]
Hence the scalar-diagonal graph series is
\begin{equation}\label{eq:ahat-generating}
 \sum_{g\ge1}F_g(\cA)x^{2g}
 =\kappa(\cA)
  \left(\frac{x/2}{\sin(x/2)}-1\right),
 \qquad
 \delta F_g^{\mathrm{cross}}(\cA)=0.
\end{equation}
The first coefficients are
\[
 \lambda_1^{\mathrm{FP}}=\frac1{24},
 \qquad
 \lambda_2^{\mathrm{FP}}=\frac7{5760},
 \qquad
 \lambda_3^{\mathrm{FP}}=\frac{31}{967680}.
\]

\begin{example}[Heisenberg realization;
\ClaimStatusConditional]
\label{ex:thmD-heis}
For the rank-one Heisenberg algebra
\([a_m,a_n]=k\,m\,\delta_{m+n,0}\mathbf1\), one has
\(\kappa(\cH_k)=k\).  The packages in Theorem~D give
\[
 \operatorname{tr}_1
 \operatorname{Obs}^{\mathrm{def}}_{1,1}(\cH_k)=k\lambda_1,
 \qquad
 \mathfrak O_g^K(\cH_k)=k\lambda_{-1}(\mathbb E_g),\quad g\ge2,
 \qquad
 \operatorname{obs}^{\mathrm{Hdg}}_g(\cH_k)=(-1)^gk\lambda_g,\quad g\ge2.
\]
On the scalar-diagonal graph lane,
\[
 \sum_{g\ge1}F_g(\cH_k)x^{2g}
 =k\left(\frac{x/2}{\sin(x/2)}-1\right).
\]
\end{example}

%%% ================================================================
%%% SECTION 7: THEOREM H
%%% ================================================================

\section{Theorem H: chiral Hochschild cohomology}
\label{sec:theorem-H}

\subsection{Statement}

Theorem~D separates the native deformation class, its Hodge shadow,
and the numerical graph trace.  Theorem~H determines the
family-specific cohomological support of the level-$3$ chart centre.

\begin{theorem}[Theorem H: family-indexed chiral Hochschild support;
\ClaimStatusConditional]
\label{thm:hochschild}
\textsc{Type signature.} Quadrant: Open. Presentation:
\[
 C^\bullet_{\mathrm{ch}}(A_b,A_b)
 :=R\!\operatorname{Hom}_{A_b\otimes A_b^{\mathrm{op}}}(A_b,A_b),
 \qquad
 \mathrm{ChirHoch}^n(A_b):=H^nC^\bullet_{\mathrm{ch}}(A_b,A_b),
\]
formed in the completed chiral bimodule category. Level: $3$
(derived chiral centre $Z^{\mathrm{der}}_{\mathrm{ch}}(A_b)$,
the algebraic closed-sector object in the Beilinson tower).
Hypothesis package: the family datum $H_H(\cA;S)$.

Fix a finite set $S\subset\mathbb Z$.  The datum
$H_H(\cA;S)$ consists of a complete filtered cochain model $Q_\cA$,
a quasi-isomorphism
\[
 \gamma_\cA\colon Q_\cA\xrightarrow{\ \sim\ }
 C^\bullet_{\mathrm{ch}}(\cA,\cA),
\]
a complete complex $K_{\cA,S}$ whose cochain terms occur in degrees
belonging to~$S$, and filtration-continuous maps
\[
 K_{\cA,S}\mathop{\xrightarrow{\ \iota_\cA\ }}Q_\cA
 \mathop{\xrightarrow{\ p_\cA\ }}K_{\cA,S},
 \qquad h_\cA\colon Q_\cA\longrightarrow Q_\cA[-1],
\]
satisfying
\[
 p_\cA\iota_\cA=\id_{K_{\cA,S}},
 \qquad
 dh_\cA+h_\cA d=\id_{Q_\cA}-\iota_\cA p_\cA.
\]
Whenever a bar model is used, the datum includes the bar--bimodule
comparison
\[
 \Gamma_{\bar B}\colon
 \operatorname{Conv}(B_X(\cA),\cA)
 \xrightarrow{\ \sim\ }C^\bullet_{\mathrm{ch}}(\cA,\cA).
\]
A cyclic pairing may further identify the shifted deformation complex
with cyclic coderivations of the completed bar; this is a separate
comparison with its suspension recorded explicitly.  The datum also
supplies convergence of the filtration and, whenever
$Q_\cA$ is first formed on ordered configurations, the
ordered-to-symmetric comparison quasi-isomorphism.  Then
\begin{equation}\label{eq:theorem-H-family-support-standalone}
 \ChirHoch^n(\cA)\cong H^n(K_{\cA,S}),
 \qquad
 \operatorname{Supp}\ChirHoch^\bullet(\cA)\subseteq S.
\end{equation}
If the displayed cohomology groups are finite-dimensional, their
Hilbert polynomial is
\begin{equation}\label{eq:hilbert-poly}
P_\cA(t)
 =\sum_{n\in S}\dim H^n(K_{\cA,S})t^n.
\end{equation}
Suppose additionally that $\cA^!$ carries $H_H(\cA^!;S^!)$ and that
$(\cA,\cA^!)$ carries a perfect degree-$d$ chain pairing.  This
separate duality datum gives
\[
 \ChirHoch^n(\cA)
 \cong
 \ChirHoch^{d-n}(\cA^!)^\vee\otimes\omega_X,
 \qquad
 P_\cA(t)=t^dP_{\cA^!}(t^{-1}).
\]
The geometric pairing determines~$d$.
\end{theorem}

\begin{remark}[Three Hochschild theories: chiral, topological, categorical]
\label{warn:three-hochschild}
Three Hochschild theories occupy distinct ambients: chiral,
topological, and categorical.
\emph{Chiral Hochschild}
$\ChirHoch^*(\cA)=
H^*R\!\operatorname{Hom}_{\cA\otimes\cA^{\mathrm{op}}}(\cA,\cA)$
lives in the completed chiral bimodule category on~$X$, and
$H_H(\cA;S)$ computes its family-specific support.
For an $E_1$ ring spectrum $R$, topological Hochschild homology is
$\mathrm{THH}(R)=R\otimes^L_{R^e}R$, equivalently factorization
homology on $S^1$, while topological Hochschild cohomology is
$F_{R^e}(R,R)$.  \emph{Categorical Hochschild}
$\mathrm{HH}^*(\mathrm{Rep}(\cA))$ governs the
Drinfeld center and the $E_2$ structure.  Comparison functors between
these theories constitute additional mathematical data.
\end{remark}

\subsection{Family-specific specializations}

\begin{proposition}[Two bounded-cohomology benchmarks;
\ClaimStatusConditional]
\label{prop:hilbert-families}
Assume that the bounded-to-chart map for the stated vertex algebra is
a quasi-isomorphism.  The computations of Bakalov--De Sole--Kac then
give the following chart cohomology.
\begin{enumerate}[label=\textup{(\roman*)}]
\item For the rank-one even free superboson $B_{\mathfrak h}$,
\[
 \bigl(\dim\ChirHoch^n(\cA_{B_{\mathfrak h}})\bigr)_{n\geq0}
 =(2,1,0,0,\ldots).
\]
\item For the Virasoro vertex algebra $\operatorname{Vir}_c$,
\[
 \dim\ChirHoch^n(\cA_{\operatorname{Vir}_c})
 =\begin{cases}
 1,&n\in\{0,2,3\},\\
 0,&n\in\mathbb Z_{\geq0}\setminus\{0,2,3\}.
 \end{cases}
\]
\end{enumerate}
These vectors are the bounded calculations of
\cite[Theorems~7.4 and~7.2]{BDSK21}; the assumed comparison carries
them to the chart complex.  Affine, principal $\cW$, lattice, and
critical-level families acquire their support from their respective
data $H_H(\cA;S_\cA)$.  At affine critical level the
Feigin--Frenkel centre supplies a degree-zero subspace modeled by
$\operatorname{Fun}(\operatorname{Op}_{\fg^\vee}(D))$.  The complete
critical-level datum $H_H$ must additionally determine every other
cohomological degree in the completed chart model.
\end{proposition}

\subsection{Proof architecture}

The identities in the datum $H_H(\cA;S)$ exhibit
$K_{\cA,S}$ and $Q_\cA$ as homotopy-equivalent cochain complexes.
Composition with $\gamma_\cA$ proves
\eqref{eq:theorem-H-family-support-standalone}.  A finite-window
construction establishes these identities stratum by stratum on the
Fulton--MacPherson compactification, verifies compatibility with the
bar faces and collision incidences, passes through the
ordered-to-symmetric comparison, and then takes the derived inverse
limit.  PBW filtrations and quadratic Koszul complexes furnish one
construction method; the resulting homotopy and convergence
identities carry the theorem.

The perfect chain pairing with the dual chart is an additional
geometric structure.  Its degree gives the reflection
$n\leftrightarrow d-n$, and finite-dimensionality then gives the
palindromic Hilbert identity.


%%% ================================================================
%%% SECTION 8: THE DEPENDENCY GRAPH
%%% ================================================================

\section{The dependency graph among Theorems A, B, C, D, and H}
\label{sec:forced-chain}

The five theorems form a typed dependency graph.  The graph has two
principal branches: reconstruction and duality at levels \(1\)--\(3\),
and modular trace at levels \(4\)--\(5\).

\[
\begin{tikzcd}[column sep=large,row sep=large]
 \cA^{\mathrm i}
   \arrow[r,"q_\cA"]
 & B_X(\cA)
   \arrow[r,"\Omega_X"]
   \arrow[d,"\mathbb D_{\Ran}"']
 & \cA
   \arrow[d,"Z^{\mathrm{der}}_{\mathrm{ch}}"] \\
 & K_X(\cA)
 & C^\bullet_{\mathrm{ch}}(\cA,\cA) \\
 K_{\cA,S}
   \arrow[r,shift left=.6ex,"\iota_\cA"]
 & Q_\cA
   \arrow[l,shift left=.6ex,"p_\cA"]
   \arrow[r,"\gamma_\cA"]
 & C^\bullet_{\mathrm{ch}}(\cA,\cA)
\end{tikzcd}
\]

The upper two rows display quadratic compression, reconstruction,
Verdier duality, and the derived centre.  The third row is the strong
deformation retract and chart comparison of Theorem~H.  Theorem~D has
the separate typed diagram
\[
\begin{tikzcd}[column sep=large,row sep=large]
 \operatorname{Obs}^{\mathrm{def}}_g(\cA)
   \arrow[r,"H_D^{\mathrm{tr}}"]
   \arrow[dr,"H_D^{\mathrm{graph}}"']
 & \mathfrak O_g^K(\cA)
   \arrow[r,"\operatorname{ch}_g"]
 & \operatorname{obs}^{\mathrm{Hdg}}_g(\cA) \\
 & F_g(\cA) &
\end{tikzcd}
\]

\begin{enumerate}[label=\textup{(\roman*)}]
\item The Ran bar--cobar equivalence of Theorem~A constructs
\(R_X(\cA)=\Omega_XB_X(\cA)\simeq\cA\).  The Verdier package
constructs the separate algebra \(K_X(\cA)=\mathbb D_{\Ran}B_X(\cA)\).

\item Theorem~B compares the quadratic coalgebra
\(\cA^{\mathrm i}\) with the full bar through
\(q_\cA\colon\cA^{\mathrm i}\to B_X(\cA)\).  Its input is
\(H_{\mathrm{CL}}\), and its conclusion is the equivalence of the
three intrinsic twisting conditions.  The fourth, spectral-sequence
condition joins them under \(H_{\mathrm{PBW}}^{\mathrm{det}}\).

\item Theorem~C organizes the five scalar measurements and their
family strata.  A represented Verdier comparison supplies any
relation between the \(\cA\)- and \(\cA^!\)-rows.

\item Theorem~D starts from the cyclic modular trace.  Its packages
compare the native deformation class with a virtual \(K\)-object,
take the signed Hodge projection, and evaluate the stable-graph
functional.

\item Theorem~H starts from the complete chiral Hochschild chart
complex.  The family datum \(H_H(\cA;S)\) gives a strong deformation
retract onto a complex supported in~\(S\).
\end{enumerate}

\begin{remark}[Primitive anchor and comparison arrows]
\label{rem:non-circular}
The full bar differential and its Maurer--Cartan element are defined
before the five theorem packages.  Theorem~A applies the Ran
bar--cobar equivalence.  Theorem~B adds a quadratic model.
Theorem~C adds family-indexed scalar measurements.  Theorem~D adds
modular trace and Hodge realization.  Theorem~H adds a complete
family support model.  Each comparison arrow therefore introduces a
new mathematical datum with a specified source and target.
\end{remark}

\subsection{One new input per theorem}

\begin{center}
\small
\renewcommand{\arraystretch}{1.15}
\begin{tabular}{llll}
\toprule
\textbf{Theorem} & \textbf{Output} & \textbf{Input}
 & \textbf{Status} \\
\midrule
A & \(R_X(\cA)\simeq\cA\), \(K_X(\cA)\)
 & Ran ambient; \(H_{\mathrm{fact}},H_{\mathrm{conv}},H_{\mathrm{VD}}\)
 & published / conjectural / conditional \\
B & \(q_\cA\colon\cA^{\mathrm i}\xrightarrow{\sim}B_X(\cA)\)
 & \(H_{\mathrm{CL}},H_{\mathrm{PBW}}^{\mathrm{det}}\)
 & conditional \\
C & five family measurements
 & represented comparison data
 & family-indexed \\
D & \(\operatorname{Obs}^{\mathrm{def}}_g,\mathfrak O_g^K,
       \operatorname{obs}^{\mathrm{Hdg}}_g,F_g\)
 & \(H_D^1,H_D^K,H_D^{\mathrm{tr}},H_D^{\mathrm{graph}}\)
 & conditional \\
H & \(\operatorname{Supp}\ChirHoch^\bullet(\cA)\subseteq S\)
 & \(H_H(\cA;S)\)
 & conditional \\
\bottomrule
\end{tabular}
\end{center}

%%% ================================================================
%%% SECTION 9: THE HEISENBERG
%%% ================================================================

\section{The Heisenberg chart: the first complete local calculation}
\label{sec:heisenberg}

\subsection{Exact local algebra}

Let \(\cH_k\) be the rank-one Heisenberg vertex algebra generated by a
weight-one field \(J\).  Its defining operator product and
\(\lambda\)-bracket are
\begin{equation}\label{eq:heis-ope}
 J(z)J(w)\sim\frac{k}{(z-w)^2},
 \qquad
 [J_\lambda J]=k\lambda .
\end{equation}
Equivalently,
\[
 [J_m,J_n]=k\,m\,\delta_{m+n,0}\mathbf1 .
\]
The normalized scalar collision seed is therefore
\begin{equation}\label{eq:heis-rmatrix}
 R_2^{\mathrm{sc}}(\cH_k)=\kappa(\cH_k)=k.
\end{equation}
This equality is the exact local computation.  The strong-generator
space is one-dimensional, while the state space contains the full
descendant tower.  At \(k=0\), the vacuum projection is an ordinary
augmentation.  At \(k\neq0\), the central term \(k\lambda\) makes the
vacuum projection a curved splitting.

\subsection{Theorem A: full Ran reconstruction}

The curved second-kind chiral bar coalgebra is
\begin{equation}\label{eq:heis-bar}
 B_X^{\mathrm{II}}(\cH_k)
 =\left(T_X^c(s^{-1}\overline{\cH}_k),
   d_B,h_k\right),
 \qquad
 h_k\colon B_X^{\mathrm{II}}(\cH_k)\longrightarrow\omega_X[2],
\end{equation}
where \(\overline{\cH}_k\) is the chosen vacuum complement and
includes the descendants of \(J\), and
\(d_B=d_{\mathrm{int}}+d_{\mathrm{coll}}\) is assembled from all
collision strata.  The curvature functional \(h_k\) is supported on
length two and represents the central collision
\begin{equation}\label{eq:heis-bar-diff}
 [J_\lambda J]=k\lambda .
\end{equation}
In the sign convention of this chapter,
\[
 d_B^2=(h_k\otimes\id-\id\otimes h_k)\Delta.
\]
Its curved cobar algebra
\(\Omega_X^{\mathrm{II}}B_X^{\mathrm{II}}(\cH_k)\) carries
\(m_0=-k\omega\).  For \(k=0\), the curvature functional vanishes and
the Ran theorem of
Francis--Gaitsgory gives
\[
 R_X(\cH_0)=\Omega_XB_X(\cH_0)\xrightarrow{\sim}\cH_0
\]
in the pro-nilpotent Ran ambient.  For \(k\neq0\), a completed curved
bar--cobar comparison belongs to the second-kind package
\(H_{\mathrm{curv}}(\cH_k)\); its conclusion is a coderived
equivalence represented by the curved counit.

Under the curved Verdier package, duality gives the separate algebra
\begin{equation}\label{eq:heis-dual}
 K_X^{\mathrm{II}}(\cH_k)
 :=\mathbb D_{\Ran}B_X^{\mathrm{II}}(\cH_k).
\end{equation}
A strict presentation of \(K_X^{\mathrm{II}}(\cH_k)\) as a curved Heisenberg dual
requires the quadratic comparison, local finite duality, and a
formality map.

\subsection{Theorem B: the finite-window comparison}

Let \(\cH_k^{\mathrm i,II}\) be the curved quadratic coalgebra
determined by the strong generator, the bracket, and the induced
length-two curvature functional \(h_k^{\mathrm i}\).  The Heisenberg
instance of Theorem~B is the curvature-preserving map
\begin{equation}\label{eq:heis-cobar}
 q_{\cH_k}^{\mathrm{II}}\colon
 \cH_k^{\mathrm i,II}\longrightarrow B_X^{\mathrm{II}}(\cH_k).
\end{equation}
The datum
\(H_{\mathrm{CL}}^{\mathrm{II}}
(\cH_k,\cH_k^{\mathrm i,II},\tau)\) consists
of explicit contractions on every finite conformal-weight window,
compatibility with collision incidences and bar faces, and the
Mittag--Leffler passage to the completed chart.  Once this datum is
constructed, it gives the curved quadratic comparison and the
second-kind cobar counit.  At \(k=0\), these specialize to the
ordinary quasi-isomorphisms of Theorem~B.  Thus the one-line OPE supplies the local seed,
and the finite-window contraction supplies the global chiral theorem.

\subsection{Theorems C and D: scalar and Hodge realizations}

The row-C input furnished directly by the OPE is
\(\kappa(\cH_k)=k\).  A represented curved Verdier trace comparison
may identify the dual scalar with \(-k\); on that comparison surface,
\begin{equation}\label{eq:heis-complementarity}
 \kappa(\cH_k)+\kappa(K_X^{\mathrm{II}}(\cH_k))=0.
\end{equation}
This is a relation between represented trace functionals.

Under the packages of Theorem~D,
\[
 \operatorname{tr}_1
 \operatorname{Obs}^{\mathrm{def}}_{1,1}(\cH_k)=k\lambda_1,
 \qquad
 \mathfrak O_g^K(\cH_k)=k\lambda_{-1}(\mathbb E_g),\quad g\ge2,
 \qquad
 \operatorname{obs}^{\mathrm{Hdg}}_g(\cH_k)=(-1)^gk\lambda_g,\quad g\ge2.
\]
On the scalar-diagonal stable-graph lane,
\[
 F_g(\cH_k)=k\lambda_g^{\mathrm{FP}},
 \qquad
 \sum_{g\ge1}F_g(\cH_k)x^{2g}
 =k\left(\frac{x/2}{\sin(x/2)}-1\right).
\]

\subsection{Theorem H: bounded benchmark and chart transport}

For the rank-one even free superboson \(B_{\mathfrak h}\),
Bakalov--De Sole--Kac prove
\[
 \dim H^n_{\mathrm{ch},b}(B_{\mathfrak h})
 =
 \begin{cases}
 2,&n=0,\\
 1,&n=1,\\
 0,&n\ge2,
 \end{cases}
 \qquad\text{\cite[Theorem~7.4]{BDSK21}}.
\]
Fix \(k\in\mathbf C^\times\) and choose \(u\in\mathbf C^\times\) with
\(u^2=k\).  The rescaling
\(\rho_u\colon B_{\mathfrak h}\xrightarrow{\sim}\cH_k\),
\(a\mapsto u^{-1}J\), transports the normalized bracket to
\([J_\lambda J]=k\lambda\).  A bounded-to-chart quasi-isomorphism
\[
 \chi^{\mathrm{bd}}_{\cH_k}\colon
 C^\bullet_{\mathrm{ch},b}(B_{\mathfrak h})
 \xrightarrow{\ C^\bullet(\rho_u)\ }
 C^\bullet_{\mathrm{ch},b}(\cH_k)
 \xrightarrow{\sim}
 C^\bullet_{\mathrm{ch}}(\cA_{\cH_k},\cA_{\cH_k})
\]
transports this vector to the chart complex.  In the language of
Theorem~H, that comparison supplies the family support datum with
\(S_{\cH}=\{0,1\}\) and cohomology vector \((2,1,0,\ldots)\).
The fiber \(k=0\) is the commutative current algebra and carries its
own family datum.

\begin{remark}[Abelian Chern--Simons comparison]
\label{rem:abelian-cs}
For the standard chiral boundary condition of abelian
Chern--Simons theory at level~\(k\), the boundary current algebra is
\(\cH_k\).  The equality \(\kappa(\cH_k)=k\) matches the OPE
coefficient with the gauge-theory level.  A comparison between the
stable-graph functional above and a perturbative
Chern--Simons amplitude is an additional open--closed map.
\end{remark}

\subsection{Summary of the worked case}

\begin{center}
\small
\renewcommand{\arraystretch}{1.15}
\begin{tabular}{lll}
\toprule
\textbf{Theorem} & \textbf{Heisenberg statement}
 & \textbf{Load-bearing datum} \\
\midrule
A & ordinary Ran counit at \(k=0\); curved coderived counit at \(k\ne0\)
 & Ran ambient; \(H_{\mathrm{curv}},H_{\mathrm{conv}},
      H_{\mathrm{VD}}^{\mathrm{II}}\) \\
B & \(q_{\cH_k}^{\mathrm{II}}\colon
      \cH_k^{\mathrm i,II}\to B_X^{\mathrm{II}}(\cH_k)\)
 & finite-window \(H_{\mathrm{CL}}^{\mathrm{II}}\) \\
C & \(\kappa(\cH_k)=k\)
 & represented dual trace for the paired row \\
D & \(\operatorname{obs}^{\mathrm{Hdg}}_g=(-1)^gk\lambda_g\)
 & \(H_D^K\) \\
H & bounded vector \((2,1,0,\ldots)\)
 & bounded-to-chart comparison \\
\bottomrule
\end{tabular}
\end{center}

The exact local input is the bracket \([J_\lambda J]=k\lambda\).
Every global conclusion records the map that carries this local input
through configurations, completion, duality, Hodge realization, or
cohomological comparison.

%%% ================================================================
%%% SECTION 10: THE STANDARD LANDSCAPE
%%% ================================================================

\section{The standard-family atlas: source data and comparison status}
\label{sec:landscape}

The standard-family atlas is finite.  Each row begins with an
exact local calculation and acquires its global entries through named
comparison maps.  The atlas is the finite census of constructed and
conditional rows.

\subsection{Local calculations before global transport}

\begin{table}[ht]
\centering
\caption{Source data, global comparison, and epistemic status.}
\label{tab:census}
\renewcommand{\arraystretch}{1.25}
\small
\begin{tabular}{@{}p{0.14\textwidth}p{0.20\textwidth}
p{0.20\textwidth}p{0.22\textwidth}p{0.14\textwidth}@{}}
\toprule
\textbf{Family} & \textbf{Local datum}
& \textbf{Benchmark} & \textbf{Global comparison}
& \textbf{Status} \\
\midrule
Heisenberg \(\cH_k\)
& \([J_\lambda J]=k\lambda\), \(\kappa=k\)
& bounded superboson vector \((2,1,0,\ldots)\)
& \(H_{\mathrm{CL}}\), \(H_D\), bounded-to-chart map
& OPE and \(\kappa\) proved; transport conditional \\

Affine \(V_k(\fg)\)
& \([J^a_\lambda J^b]=[J^a,J^b]+k\lambda(a,b)\)
& \(\displaystyle
 \kappa=\frac{(k+h^\vee)\dim\fg}{2h^\vee}\)
& quadratic contraction, represented trace, family \(H_H\)
& \(\kappa\) proved in trace normalization; transport conditional \\

\(\beta\gamma_\lambda\), \(bc_\lambda\)
& free-field Wick algebra
& \(\kappa_{\beta\gamma}=6\lambda^2-6\lambda+1\),
  \(\kappa_{bc}=-\kappa_{\beta\gamma}\)
& collision-residue descent and family support datum
& Wick values computed; transport conditional \\

Virasoro \(\operatorname{Vir}_c\)
& \(T(z)T(w)\) OPE, \(\kappa=c/2\)
& level-four Gram coefficient
  \(\displaystyle 10/[c(5c+22)]\), \(c(5c+22)\ne0\);
  bounded support \(\{0,2,3\}\)
& \(H_{\mathrm{res}}\), bounded-to-chart map, \(H_D\)
& Gram seed exact on domain; chart support conditional \\

Principal \(\cW_N\)
& strong weights \(2,\ldots,N\)
& \(\displaystyle\kappa=c(H_N-1)\)
& mixed-channel graph functional and family \(H_H\)
& value conditional on DS/bar trace comparison \\

\(K3\times E\) comparison row
& \(e_{\mathrm{top}}(K3)=24\), \(\mathrm{wt}(\Delta_5)=5\)
& native \((0,0,\star,5,24)\); Open entry~$3$
& realization package \(H_{\mathsf B}\) across Open and CY quadrants
& four native entries exact; Open entry conditional \\
\bottomrule
\end{tabular}
\end{table}

The local-datum and benchmark columns record calculations in their
stated normalizations.  The comparison column names the map into the
chiral bar, Hodge, Hochschild, or cross-volume ambient; the final
column records the resulting epistemic status.

\subsection{The census-scoped archetypes}

The atlas uses the labels
\[
 \mathsf G,\quad\mathsf L,\quad\mathsf C,\quad
 \mathsf M,\quad\mathsf B.
\]
Rows \(\mathsf G,\mathsf L,\mathsf C,\mathsf M\) refer to the
finite Vol~I standard-family census.  Row \(\mathsf B\) is the
conditional \(K3\times E\) comparison row.  Their selected
presentation depths
\[
 2,\quad3,\quad4,\quad\infty,\quad5
\]
belong to the selected shadow presentations.  The Vol~I scalar depth
and the Vol~III Borcherds comparison depth are separate constructions.
A family receives an archetype label after its residue projection and
recurrence data have been constructed.

\begin{remark}[Five measurements]
The five coordinates
\[
 \kappa_{\mathrm{cat}},\quad
 \kappa_{\mathrm{ch}}^{\mathrm{Hodge}},\quad
 \kappa_{\mathrm{ch}}^{\mathrm{Heis}},\quad
 \kappa_{\mathrm{BKM}},\quad
 \kappa_{\mathrm{fiber}}
\]
have distinct sources and targets.  A row in the \(5\times5\) table
is a tuple of comparisons among these measurements.  The
\(K3\times E\) coordinate records a categorical value, two chiral
values, a Borcherds weight, and a topological Euler characteristic.
\end{remark}

\subsection{The Virasoro scalar lane}

For \(c(5c+22)\ne0\), the exact local seeds are
\[
 R_2=\frac c2,\qquad R_3=2,\qquad
 R_4^{\mathrm{Gram}}=\frac{10}{c(5c+22)}.
\]
On the domain \(c(5c+22)\ne0\), the weighted unordered-pair
recurrence gives
\[
 R_5^{\mathrm{Ricc}}
 =-\frac{48}{c^2(5c+22)}
\]
and the sixth weighted-Riccati coefficient
\[
 R_6^{\mathrm{Ricc}}
 =\frac{80(45c+193)}
 {3c^3(5c+22)^2}.
\]
Imposing the formal quadratic relation
\[
 2R_2C_6^{\mathrm{rel}}+2R_3R_5+R_4^2=0
\]
with \(R_2=c/2\), \(R_3=2\),
\(R_4=10/[c(5c+22)]\), and
\(R_5=-48/[c^2(5c+22)]\), determines the order-six relation coefficient
\[
 C_6^{\mathrm{rel}}
 =\frac{4(240c+1031)}
 {c^3(5c+22)^2}.
\]
A singular-vector interpretation is supplied by an explicit level-six
radical vector and its decoupling map to the chosen quotient or residue
complex.
A datum \(H_{\mathrm{res}}(\operatorname{Vir}_c;X)\) consists of an
ordered residue complex, an Arnold class, a scalar projection, and
its normalization.  This datum compares connected Ward correlators,
Gram coefficients, the weighted recurrence, and the formal relation
coefficient in a common scalar ambient.

For \(45c+218\ne0\), the discriminant of the weighted generating
series has a simple zero and hence a square-root branch point.  Its
coefficients have the corresponding geometric growth with the
power-law correction supplied by singularity analysis.

\subsection{Family-indexed Hochschild support}

For every chart \(\cA\), Theorem~H begins with a complete complex
\(Q_\cA\), a support model \(K_{\cA,S}\), and a strong deformation
retract.  The resulting Hilbert polynomial, when the cohomology is
finite-dimensional, is
\[
 P_\cA(t)=\sum_{n\in S}\dim H^n(K_{\cA,S})t^n.
\]
Two bounded calculations provide exact benchmarks:
\[
 \bigl(\dim H^n_{\mathrm{ch},b}(B_{\mathfrak h})\bigr)_{n\ge0}
 =(2,1,0,\ldots),
 \qquad
 \operatorname{Supp}
 H^\bullet_{\mathrm{ch},b}(\operatorname{Vir}_c)=\{0,2,3\},
\]
with one-dimensional Virasoro groups in those three degrees
\cite[Theorems~7.4 and~7.2]{BDSK21}.  Each chart-level use of these
vectors names its bounded-to-chart quasi-isomorphism.  Affine,
principal \(\cW\), lattice, free-field, and critical-level rows carry
their respective data \(H_H(\cA;S_\cA)\).

\subsection{Provenance rule for every row}

A numerical row is assembled in the following order:
\[
 \text{OPE/root/character calculation}
 \longrightarrow
 \text{chain model}
 \longrightarrow
 \text{comparison quasi-isomorphism}
 \longrightarrow
 \text{global invariant}.
\]
The source calculation carries status \emph{computed} or
\emph{proved elsewhere}.  The comparison carries status
\emph{proved}, \emph{conditional}, or \emph{open}.  The global row
inherits the status of its comparison.  This ordering makes the
landscape a reproducible atlas of mathematical constructions.

\begin{thebibliography}{99}

\bibitem{Arnold69}
V.~I.~Arnold,
The cohomology ring of the colored braid group,
\emph{Mat. Zametki} \textbf{5} (1969), 227--231.

\bibitem{BD04}
A.~Beilinson and V.~Drinfeld,
\emph{Chiral Algebras},
AMS Colloquium Publications, vol.~51, 2004.

\bibitem{BGS96}
A.~Beilinson, V.~Ginzburg, and W.~Soergel,
Koszul duality patterns in representation theory,
\emph{J.\ Amer.\ Math.\ Soc.}\ \textbf{9} (1996), no.~2, 473--527.

\bibitem{BDSK21}
B.~Bakalov, A.~De Sole, and V.~G.~Kac,
\emph{Computation of cohomology of vertex algebras},
Japan. J. Math. \textbf{16} (2021), 81--154,
arXiv:2002.03612.

\bibitem{Borcherds1998}
R.~E.~Borcherds,
Automorphic forms with singularities on Grassmannians,
\emph{Invent. Math.} \textbf{132} (1998), no.~3, 491--562.

\bibitem{CG17}
K.~Costello and O.~Gwilliam,
\emph{Factorization Algebras in Quantum Field Theory}, vols.~1--2,
Cambridge University Press, 2017--2021.

\bibitem{FG12}
J.~Francis and D.~Gaitsgory,
Chiral Koszul duality,
\emph{Selecta Math.\ (N.S.)}\ \textbf{18} (2012), no.~1, 27--87.

\bibitem{FM94}
W.~Fulton and R.~MacPherson,
A compactification of configuration spaces,
\emph{Ann. of Math.\ (2)} \textbf{139} (1994), no.~1, 183--225.

\bibitem{FP03}
C.~Faber and R.~Pandharipande,
Hodge integrals, partition matrices, and the $\lambda_g$ conjecture,
\emph{Ann. of Math.\ (2)} \textbf{157} (2003), no.~1, 97--124.

\bibitem{FrenkelBenZvi}
E.~Frenkel and D.~Ben-Zvi,
\emph{Vertex Algebras and Algebraic Curves},
second edition, Mathematical Surveys and Monographs, vol.~88,
American Mathematical Society, 2004.

\bibitem{GarlandLepowsky}
H.~Garland and J.~Lepowsky,
Lie algebra homology and the Macdonald--Kac formulas,
\emph{Invent.\ Math.}\ \textbf{34} (1976), 37--76.

\bibitem{GetzlerKapranov98}
E.~Getzler and M.~M.~Kapranov,
Modular operads,
\emph{Compositio Math.}\ \textbf{110} (1998), 65--126.

\bibitem{GLZ24}
J.~Gui, H.~Li, and J.~Zeng,
Quadratic duality for chiral algebras,
\emph{Adv.\ Math.}\ \textbf{451} (2024), Paper no.~109791.

\bibitem{Gritsenko1999}
V.~Gritsenko and V.~Nikulin,
The Igusa modular forms and ``the simplest'' Lorentzian
Kac--Moody algebras,
\emph{Mat. Sb.} \textbf{187} (1996), no.~11, 27--66;
arXiv:alg-geom/9603010.

\bibitem{LV12}
J.-L.~Loday and B.~Vallette,
\emph{Algebraic Operads},
Grundlehren der mathematischen Wissenschaften, vol.~346,
Springer, 2012.

\bibitem{Mok25}
S.~C.~Mok,
Logarithmic Fulton--MacPherson configuration spaces,
arXiv:2503.17563, 2025.

\bibitem{Mukai1987}
S.~Mukai,
On the moduli space of bundles on K3 surfaces, I,
in \emph{Vector Bundles on Algebraic Varieties}, Tata Institute of
Fundamental Research Studies in Mathematics, vol.~11, 1987,
341--413.

\bibitem{Mumford83}
D.~Mumford,
Towards an enumerative geometry of the moduli space of curves,
in \emph{Arithmetic and Geometry}, vol.~II, Birkh\"auser, 1983,
271--328.

\bibitem{Faltings84}
G.~Faltings,
Calculus on arithmetic surfaces,
\emph{Ann.\ of Math.\ (2)} \textbf{119} (1984), 387--424.

\bibitem{MS24}
F.~Malikov and V.~Schechtman,
Homotopy chiral algebras,
arXiv:2408.16787, 2024.

\bibitem{Priddy70}
S.~Priddy,
Koszul resolutions,
\emph{Trans.\ Amer.\ Math.\ Soc.}\ \textbf{152} (1970), 39--60.

\bibitem{Positselski11}
L.~Positselski,
\emph{Two kinds of derived categories, Koszul duality, and
comodule-contramodule correspondence},
Mem.\ Amer.\ Math.\ Soc.\ \textbf{212} (2011), no.~996.

\end{thebibliography}

%%% ================================================================
%%% ADDENDUM
%%% ================================================================

\appendix
\section{Comparison morphisms for the five theorems}
\label{sec:five-theorems-ambients}

This appendix records the mathematical object consumed by each
theorem and the construction that remains family-dependent.

\subsection*{Theorem A}

The published ambient statement is enhanced bar--cobar equivalence
for the reduced associative operad in the pro-nilpotent Ran category.
The factorization restriction is the closure package
\(H_{\mathrm{fact}}\): preservation of factorization equivalences,
collision base change, units, and completed inverse limits.
The Verdier branch is the dualizability package
\(H_{\mathrm{VD}}\).  A self-duality point additionally carries a
map \(\sigma_\cA\colon\cA\xrightarrow{\sim}K_X(\cA)\).

\subsection*{Theorem B}

The family calculation is the cone
\[
 \operatorname{Cone}\!\left(
 q_\cA\colon\cA^{\mathrm i}\longrightarrow B_X(\cA)
 \right).
\]
The package \(H_{\mathrm{CL}}\) builds this cone on finite conformal
weight windows, contracts its positive-weight part compatibly with
collision incidences and bar faces, and passes to the completed
inverse limit.  This construction proves the three intrinsic
twisting conditions of Theorem~B.  The detecting PBW package
\(H_{\mathrm{PBW}}^{\mathrm{det}}\) adds the equivalent $E_2$-collapse
and diagonal-support condition.

\subsection*{Theorem C}

A row of the \(5\times5\) matrix is a tuple of five measurements.
Every equality between two cells is represented by a comparison map.
The Open--CY row \(\mathsf B\) carries the realization package
\(H_{\mathsf B}\); its exact native entries are
\((0,0,\star,5,24)\), with
\(e_{\mathrm{top}}(K3)=24\) and \(\mathrm{wt}(\Delta_5)=5\).
The comparison \(\eta_{K3,E}\) supplies the Open Heisenberg entry~$3$.

\subsection*{Theorem D}

The hierarchy is
\[
 \operatorname{Obs}^{\mathrm{def}}_g
 \xrightarrow{\,H_D^{\mathrm{tr}}\,}
 \mathfrak O_g^K
 \xrightarrow{\,\operatorname{ch}_g\,}
 \operatorname{obs}^{\mathrm{Hdg}}_g,
 \qquad
 \operatorname{Obs}^{\mathrm{def}}_g
 \xrightarrow{\,H_D^{\mathrm{graph}}\,}
 F_g.
\]
The comparison \(H_D^K\) identifies
\(\mathfrak O_g^K\) with
\(\kappa\lambda_{-1}(\mathbb E_g)\).
The Chern-root calculation then gives
\((-1)^g\kappa\lambda_g\).  The graph functional gives
\(\kappa\lambda_g^{\mathrm{FP}}+\delta F_g^{\mathrm{cross}}\).

\subsection*{Theorem H}

The family datum \(H_H(\cA;S)\) is the strong deformation retract
\[
 K_{\cA,S}
 \mathop{\xrightarrow{\ \iota\ }}Q_\cA
 \mathop{\xrightarrow{\ p\ }}K_{\cA,S},
 \qquad
 dh+hd=\id_{Q_\cA}-\iota p.
\]
It includes the chart quasi-isomorphism, completion, and
ordered-to-symmetric descent.  Bounded vertex-cohomology calculations
enter through a separate bounded-to-chart quasi-isomorphism.

\section{Exact computational layer}
\label{sec:standalone-exact-computation}

The Virasoro Ward recursion has been evaluated exactly through arity
six.  At central charge \(c=1\) and coordinates
\((z_1,\ldots,z_n)=(0,\ldots,n-1)\), the connected values at arities
five and six are
\[
 G^{\mathrm{conn}}_5=\frac{775}{5184},
 \qquad
 G^{\mathrm{conn}}_6=\frac{49705}{373248}.
\]
Möbius inversion on set partitions and the free-boson cycle expansion
give the same rational functions.  Their passage to a scalar
\(S_r\) is the residue map \(H_{\mathrm{res}}\).

On the domain \(c(5c+22)\ne0\), the weighted unordered-pair recurrence
gives
\[
 R_5^{\mathrm{Ricc}}
 =-\frac{48}{c^2(5c+22)},\qquad
 R_6^{\mathrm{Ricc}}
 =\frac{80(45c+193)}{3c^3(5c+22)^2}.
\]
The formal relation construction gives
\[
 C_6^{\mathrm{rel}}
 =\frac{4(240c+1031)}{c^3(5c+22)^2}.
\]
These are exact outputs of two named formal constructions.

For exceptional Yangians, exact computation presently covers root
systems, Casimir eigenvalues, and tensor-square decompositions.
A Yang--Baxter evaluation requires triple-tensor projectors and
recoupling data.  For Drinfeld--Sokolov reduction, exact computation
covers type-\(A\) root arithmetic, nilpotent-orbit dimensions, and
the \(\mathfrak{sl}_2\) central-charge formulas.  The induced
chiral Hochschild map requires an explicit BRST chain map and
homotopy.

\section{Open construction problems}
\label{sec:standalone-open-obligations}

\begin{enumerate}[label=\textup{(O\arabic*)}]
\item Construct the Heisenberg finite-window collision contraction in
all strata and prove its compatibility with the inverse system.

\item Construct \(H_{\mathrm{res}}(\operatorname{Vir}_c;X)\).  Its
residue map defines \(S_6(\operatorname{Vir}_c;H_{\mathrm{res}})\) and
compares this geometric scalar with the weighted-Riccati and formal
relation coefficients.  A singular-vector comparison identifies its
null-state specializations.

\item Construct the family data \(H_H(\cA;S_\cA)\) for affine and
principal \(\cW\) charts, including their bounded-to-chart maps.

\item Construct \(\beta_g(\cA)\) in \(H_D^K\) from the scalar virtual
bar object to \(\kappa(\cA)\lambda_{-1}(\mathbb E_g)\), and construct
the mixed-channel stable-graph functional.

\item Construct the Open--CY realization package
\(H_{\mathsf B}\) for the \(K3\times E\) row and compare its five
coordinates in their native categories.

\item Construct the triple-tensor recoupling data for exceptional
Yangians and the BRST-induced chain maps for Drinfeld--Sokolov
reduction.
\end{enumerate}

Each problem names a source, a target, and the map whose construction
advances the corresponding theorem.

\end{document}
